\documentclass[12pt, a4paper]{article}
\usepackage[T1]{fontenc}
\usepackage{lmodern}
\usepackage{amsmath, amsthm, amssymb, mathtools}
\usepackage[margin=1in]{geometry}
\usepackage{xr-hyper}
\usepackage[colorlinks=true,linkcolor=blue,citecolor=blue]{hyperref}
\usepackage{cleveref}
\usepackage{graphicx, subcaption}
\usepackage{natbib}
\usepackage{booktabs}
\usepackage{enumerate}

% Code references: scaled monospace with line-breaking
\newcommand{\code}[1]{\texttt{\small #1}}

\newtheorem{theorem}{Theorem}[section]
\newtheorem{lemma}[theorem]{Lemma}
\newtheorem{proposition}[theorem]{Proposition}
\newtheorem{corollary}[theorem]{Corollary}
\newtheorem{definition}{Definition}[section]
\newtheorem{remark}{Remark}[section]
\newtheorem{conjecture}[theorem]{Conjecture}
\newtheorem{derivation}[theorem]{Derivation}
\newtheorem{observation}[theorem]{Observation}
\numberwithin{equation}{section}

\newcommand{\dd}{\mathrm{d}}
\newcommand{\seriestitle}{Why Rotating Fluids Make Polygons}


\externaldocument[II-]{../paper-2-physics/main}
\externaldocument[III-]{../paper-3-planets/main}
\externaldocument[A-]{../paper-A-appendices/main}

\title{\seriestitle\\[6pt]\large Paper I: Mathematical Foundations}
\author{Gordon Speagle}
\date{}

\begin{document}
\maketitle

\begin{abstract}
We establish the mathematical foundations for polygon formation in two-dimensional fluids: the sign rule, the Havelock identity, stability thresholds on constant-curvature surfaces, and the algebraic structure of the stability boundary. The five algebraic-integer thresholds on the hyperbolic plane are Pell units in quadratic fields, terminating at the golden ratio $\xi^*(23) = \varphi^{-2}$, and correspond to closed geodesics on the modular surface $\mathrm{SL}(2,\mathbb{Z})\backslash\mathbf{H}^2$ with traces $c = 2 + 16/(N{-}7)$. On arithmetic quotients $\Gamma_0(p)\backslash\mathbf{H}^2$, the Legendre symbol $(\Delta/p)$ determines which thresholds are topologically protected. The palindromic structure generalises to arbitrary genus: on a surface with trace field of degree~$d$, the threshold satisfies a palindromic polynomial of degree~$2d$.
\end{abstract}

\section{Introduction}
\label{sec:intro}

The stability eigenvalue of a regular $N$-gon of point vortices on
a two-dimensional surface admits a decomposition
\begin{equation}\label{eq:central-decomposition}
  \lambda_m = C_1(\text{surface}) - \frac{m(N-m)}{2},
  \qquad m = 1, \ldots, N-1,
\end{equation}
that separates a single surface-dependent number $C_1$ from a
universal function $f(m,N) = m(N{-}m)/2$ determined entirely by the
representation theory of the cyclic group $\mathbb{Z}_N$.
Every stability result on every constant-curvature surface---flat
plane, sphere, hyperbolic plane---follows from this formula with
one number substituted in.

This decomposition is not new: it is implicit in
\citet{Havelock1931} and explicit in \citet{BoattoCabral2003}.
What has not been made fully explicit is the depth of its
consequences.  The present paper develops five:
the $m$-dependent part $f(m,N)$ is metric-independent
(Section~\ref{sec:havelock-identity});
the sign rule for $H''(0) < 0$ follows from the logarithm's
dilation invariance (Section~\ref{sec:sign-rule});
the stability threshold $N_{\mathrm{crit}} = 7$ is a closed-form
output of the $\mathbb{Z}_N$ Casimir
(Corollary~\ref{cor:why-7});
the curvature-stability dichotomy
$\mathrm{sgn}(N_{\mathrm{crit}} - 7) = -\mathrm{sgn}(K)$
is topological
(Proposition~\ref{prop:curvature-dichotomy});
$N_{\mathrm{crit}} = 7$ is a topological universal,
holding on every quotient $\Gamma\backslash\mathbf{H}^2$ at
the flat-plane limit regardless of $\Gamma$
(Theorem~\ref{thm:universality});
and the threshold polynomials on $\mathbf{H}^2$ connect to Pell
equations, the golden ratio, and algebraic number theory
(Section~\ref{sec:algebraic-structure}).

The physical motivation---polygonal storm patterns on Saturn
and Jupiter \citep{Godfrey1988, Adriani2018}---is developed in
Paper~III.  The constrained energy landscape and
continuous-vorticity bridge are developed in Paper~II.
Paper~I is self-contained: it requires only the definitions of
point vortex energy and the Havelock eigenvalue, both recalled
below.

\paragraph{What this paper proves.}
\begin{enumerate}
  \item \textbf{The sign rule} (Theorem~\ref{thm:sign-rule}):
    for any pairwise interaction $h(r)$ with $h' < 0$ and
    $(rh')' \le 0$, the $N$-gon is a saddle of the unconstrained
    energy with $H''(0) < 0$ along the centered spiral deformation.
    The logarithmic interaction $h = -\ln r$ is the unique
    dilation-invariant borderline case (Section~\ref{sec:log-interaction}).
  \item \textbf{The Havelock identity and its metric-independence}
    (Section~\ref{sec:havelock-identity}):
    the $m$-dependent part $f(m,N) = m(N{-}m)/2$ of the Havelock
    eigenvalue is a consequence of $\mathbb{Z}_N$ representation
    theory alone; only the additive constant $C_1$ carries geometric
    information.
  \item \textbf{Stability thresholds on constant-curvature surfaces}
    (Sections~\ref{sec:curvature-thresholds}--\ref{sec:algebraic-structure}):
    the curvature-stability dichotomy
    (Proposition~\ref{prop:curvature-dichotomy}),
    the Morse index formula $\mu(N) = N - 5$ for $N \ge 8$
    (Proposition~\ref{prop:morse-index}), the closed-form derivation
    of $N_{\mathrm{crit}} = 7$ (Corollary~\ref{cor:why-7}),
    the monotone spectral flow from flat to hyperbolic geometry,
    and the robustness of $\mu(N)$ under mode-coupling
    perturbations (Remark~\ref{rmk:topological-protection}).
  \item \textbf{Algebraic structure of the stability boundary}
    (Section~\ref{sec:algebraic-structure}):
    the field extension pattern, the Pell equation connection,
    the algebraic-integer classification
    $N \in \{8, 9, 11, 15, 23\}$
    (Proposition~\ref{prop:alg-int}), the golden ratio terminal case
    $\xi^*(23) = \varphi^{-2}$, the continued fraction structure
    (Proposition~\ref{prop:cf-palindromic}),
    the geodesic correspondence with the modular surface
    (Proposition~\ref{prop:geodesic-correspondence}),
    the higher-genus palindromic hierarchy
    (Proposition~\ref{prop:higher-genus}),
    and the Legendre-symbol selection rule on arithmetic surfaces
    (Proposition~\ref{prop:legendre-selection}).
  \item \textbf{The RG classification}
    (Proposition~\ref{prop:rg-classification}):
    the logarithm is the unique IR fixed point of the additive
    coarse-graining map, with the sign rule demarcating its basin of
    attraction.
\end{enumerate}

\paragraph{What this paper does not prove.}
The constrained energy minimum theorem---showing the $N$-gon
\emph{minimises} interaction energy on the conserved-quantity surface
for $N \le 7$---is proved in Paper~II (Theorem~8.2), along with the
central vortex stabilization, the $N = 7$ quartic bifurcation, and the
blob bridge to continuous vorticity.
The planetary applications (Saturn hexagon, Jupiter polygons, Neptune
predictions) and the constraint-intersection selection principle are
developed in Paper~III.

\paragraph{Structure.}
Section~\ref{sec:log-interaction} establishes the logarithmic
interaction's uniqueness and symmetry properties.
Section~\ref{sec:sign-rule} proves the sign rule.
Sections~\ref{sec:havelock-identity}--\ref{sec:curvature-thresholds}
develop the Havelock identity, the Riemannian generalisation, and
the stability thresholds on curved surfaces.
Section~\ref{sec:algebraic-structure} explores the algebraic number
theory of the threshold polynomials.


%% ====================================================================

%% ====================================================================
\label{part:math}
%% ====================================================================

%% ====================================================================
\section{The logarithmic interaction: uniqueness and symmetry}
\label{sec:log-interaction}

\subsection{Dilation invariance and the Cauchy functional equation}
\label{sec:dilation-invariance}

\begin{remark}[The logarithm as the unique dilation-invariant interaction]
\label{rem:rg-fixed-point}
Among all pairwise interactions $h:(0,\infty)\to\mathbb{R}$ with
$h'<0$, the logarithmic interaction $h(r)=-\ln r$ is the
\emph{unique interaction invariant under spatial dilation up to an
additive constant}.  Define the scale map $T_b[h](r) = h(br)$
($b>0$).  Dilation-invariance means $h(br) = h(r) + c(b)$ for
all $r$---that is, rescaling changes $h$ only by an additive constant
(irrelevant to the force $-h'$).  This is a Cauchy functional equation
whose continuous solutions with $h'<0$ are exactly $h(r)=-a\ln r + b$,
$a>0$, confirming class~(ii) as the unique dilation-invariant class.
The $N$-body lift of this per-pair statement is
Corollary~\ref{cor:dilation-rg}: the M\"obius covariance of $H$
(Proposition~\ref{prop:mobius}), restricted to the dilation subgroup,
gives $H(bz_1,\dots,bz_N)=H(z_1,\dots,z_N)-\tfrac{N(N-1)}{2}\ln b$,
which is position-independent---the same fixed-point condition at the
$N$-body level.
This symmetry is not merely a geometric observation about the
logarithm: it is the structural reason for the \emph{scale-independence}
of $N_{\mathrm{crit}}=7$.  Because scaling changes $H$ by an
additive constant (Corollary~\ref{cor:dilation-rg}), the Havelock
eigenvalue $\lambda_m = (N{-}1) - m(N{-}m)/2$---which determines
whether the $N$-gon is a stable constrained minimum---is
scale-independent, and the threshold $N=7$ cannot drift under
spatial rescaling.  Prior treatments of M\"obius covariance in vortex
dynamics (e.g.\ \citealt{Newton2001}) identify it as a conformal
symmetry of the stream function; Corollary~\ref{cor:dilation-rg}
shows it additionally implies that the stability threshold is
dilation-invariant.
\end{remark}

\begin{proposition}[RG classification of perturbations]
\label{prop:rg-classification}
Write $h(r) = -\ln r + \varepsilon\,f(r)$ where $f$ is a perturbation.
The additive coarse-graining map
$T_b\colon h \mapsto h(b\,\cdot) - h(b)$ acts on~$f$ as
$(T_b f)(r) = f(br) - f(b)$.
\begin{enumerate}
  \item \emph{Irrelevant} ($f(r) \to 0$ as $r \to \infty$):
    the perturbation is suppressed at large scales;
    $h$ flows toward $-\ln r$ under coarse-graining.
  \item \emph{Relevant} ($f(r) \to \infty$ as $r \to \infty$):
    the perturbation dominates at large scales;
    $h$ flows away from $-\ln r$.
  \item \emph{Marginal} ($f(r) \to \mathrm{const}$):
    the perturbation shifts $h$ by a constant.
\end{enumerate}
Eigenfunctions $f(r) = r^\alpha$:
$(T_b\, r^\alpha)(r) = b^\alpha(r^\alpha - 1)$,
so $r^\alpha$ with $\alpha < 0$ is irrelevant and
$\alpha > 0$ is relevant.
\end{proposition}

\begin{proof}
For $f(r) = r^\alpha$ with $\alpha < 0$:
$|f(br)| = b^\alpha |r^\alpha| \to 0$ as $b \to \infty$ for
fixed $r$, so coarse-graining suppresses the perturbation.
For $\alpha > 0$: $|f(br)| = b^\alpha |r^\alpha| \to \infty$,
so the perturbation grows under coarse-graining.
\end{proof}

\noindent\emph{Connection to the sign rule.}
For $h = -\ln r + \varepsilon\,r^\alpha$:
$h' = -1/r + \varepsilon\alpha r^{\alpha-1}$;
the sign rule condition $h' < 0$ holds at all $r > 0$ iff
$\alpha \le 0$ (irrelevant perturbations).  For $\alpha > 0$ the
sign rule is violated at sufficiently large~$r$:
\emph{relevant perturbations destroy polygon stability}.
Thus the sign rule condition~(i) is equivalent to the requirement
that $h$ lies in the \emph{basin of attraction} of the logarithmic
fixed point.

\smallskip
\noindent\emph{The stable manifold from multipole expansion.}
The basin of attraction arises naturally from compact vorticity.
For $|z|<\varepsilon\ll|w|=d$, expanding
$\ln|w-z| = \ln d - \mathrm{Re}(z/w) + O(\varepsilon^2/d^2)$
and integrating against $\omega(z)$ with centroid at the origin
gives the effective interaction
$H_\mathrm{eff}(d)=-K^2\ln d + O(\varepsilon^2/d^2)$ for
$d\gg\varepsilon$.  The corrections are irrelevant ($\alpha = -2$),
so \emph{any compact vortex distribution flows to the logarithm
at large separation} (Proposition~\ref{prop:rg-classification}).

\smallskip
\noindent\emph{Unifying the blob and $K_0$ bridges.}
This RG perspective unifies two results proved later in the paper:
\begin{itemize}
  \item The blob bridge
    (the blob bridge (Paper~II)):
    $h_\varepsilon = -\tfrac{1}{2}\ln(r^2 + 2\varepsilon^2)
    = -\ln r + O(\varepsilon^2/r^2)$.  The blob correction
    is an irrelevant perturbation ($\alpha = -2$);
    the $O(\varepsilon^2)$ convergence rate of
    the blob correction (Paper~II, Proposition~12.1) is the approach rate to the
    fixed point.
  \item The $K_0$ interaction
    (the $K_0$ interaction (Paper~II)):
    $K_0(r/R_d) = -\ln(r/R_d) - \gamma_E + r^2/(4R_d^2) + O(r^4)$
    for $r \ll R_d$.
    The exponential tail is irrelevant (faster than any power);
    the deformation-radius limit $R_d \to \infty$ is the RG flow
    toward the fixed point.
    The interaction falls in class~(iii) of the sign rule
    ($h' < 0$ but $(rh')' > 0$), so the sign rule does not
    directly guarantee $H''(0) < 0$.
    However, the \emph{leading} correction to $-\ln r$ is
    $r^2/(4R_d^2)$ ($\alpha = 2$), which is isotropic-marginal
    (Proposition~\ref{prop:alpha-2-marginal} below): it shifts
    the precession rate but not the stability eigenvalues.
    The first stability-relevant correction is $O(r^4/R_d^4)$
    ($\alpha = 4$), so the barotropic $N_{\mathrm{crit}} = 7$
    persists to $O(\eta^2)$ in the ratio $\eta = R/R_d$.
    There is no closed-form $C_1(\eta)$ for $K_0$---the Havelock
    sums involve Bessel functions at the chord distances
    $d_p = 2R\sin(\pi p/N)$, which do not simplify
    algebraically---but the numerical thresholds are given in
    Paper~II, \S13.
\end{itemize}
In the Kraichnan inverse cascade~\citep{Kraichnan1967}, vorticity
concentrates into compact structures of growing scale $L(t)$, driving
the inter-blob interaction toward the logarithm with corrections
suppressed as $(L/d)^2 \to 0$.  The cascade is the physical
RG flow; the sign rule demarcates its basin of attraction.

\begin{proposition}[Isotropic marginality at $\alpha = 2$]
\label{prop:alpha-2-marginal}
For $h(r) = -\ln r + \varepsilon\,r^2$, the Havelock eigenvalue
$\lambda_m$ is independent of $\varepsilon$ at every binding mode
$m = 2, \ldots, N{-}2$.
\end{proposition}

\begin{proof}
The pair Hessian of $g(r) = r^\alpha$ decomposes as
\begin{equation}\label{eq:anisotropy}
  \frac{\partial^2 g}{\partial x_i\,\partial x_j}\bigg|_{|x|=d}
  = \underbrace{\alpha(\alpha{-}2)\,d^{\alpha-2}}_{\displaystyle A(\alpha,d)}\,
  \frac{x_i x_j}{d^2}
  \;+\;
  \underbrace{\alpha\,d^{\alpha-2}}_{\displaystyle B(\alpha,d)}\,
  \delta_{ij},
\end{equation}
where $A$ is the \emph{radial anisotropy} and $B$ is the
\emph{isotropic part}.
\emph{Step 1}: at $\alpha = 2$, $A(2,d) = 2(2{-}2)\,d^0 = 0$.
The Hessian perturbation $\delta\nabla^2 H$ is therefore
$B\,I = 2\,I$ at every pair, making the off-diagonal blocks
of the full $2N\times 2N$ Hessian proportional to the identity.
\emph{Step 2 (Fourier block).}
Since the off-diagonal blocks are proportional to~$I_2$, the
$xx$ and $yy$ Fourier components are equal and the $xy$ component
vanishes; each $2 \times 2$ Fourier block is scalar.
For the $xx$ component at mode $m \ge 2$: the diagonal block
equals $-\sum_{p=1}^{N-1} B = -2(N{-}1)$, and the off-diagonal
sum is $\sum_p B\,e^{-2\pi i pm/N} = 2\sum_p e^{-2\pi i pm/N} = -2$
(since $\sum_{p=1}^{N-1} e^{-2\pi i pm/N} = -1$ for $m \ge 1$).
The Fourier block is therefore $-2(N{-}1) + (-2) = -2N$,
identically for $xx$ and $yy$.  The Rayleigh quotient of the
radial Fourier mode $v_m$ gives
$v_m^\top\,\delta\nabla^2\!H\,v_m = -2N$.

\emph{Step 3 (rotation frequency).}
The equilibrium condition for $H = -\sum G(z_j,z_k)$ at the
$N$-gon gives $\delta\Omega = -(1/R)\,\partial(\delta H)/\partial x_0$,
where $\delta H = -\sum g(d_{jk})$.
For $g(r) = r^2$: $g'(r) = 2r$, and the radial gradient at
vortex $0 = (R,0)$ is
$\partial(\delta H)/\partial x_0 =
-\sum_{p=1}^{N-1} g'(d_p)\,(x_0{-}x_p)/d_p
= -\sum_{p=1}^{N-1} 2d_p\cdot(x_0{-}x_p)/d_p
= -2\sum_{p=1}^{N-1}(x_0 - x_p)$.
Now $\sum_{p=0}^{N-1} x_p = R\sum_{p=0}^{N-1}\cos(2\pi p/N) = 0$
(centroid at the origin), so
$\sum_{p=1}^{N-1}(x_0 - x_p) = (N{-}1)x_0 - \sum_{p=1}^{N-1}x_p
= (N{-}1)R - (-R) = NR$.
Hence $\partial(\delta H)/\partial x_0 = -2NR$ and
$\delta\Omega = -(-2NR)/R = 2N$.

\emph{Step 4 (cancellation).}
$\delta\lambda_m = -\delta\Omega - v_m^\top\,\delta\nabla^2\!H\,v_m
= -(2N) - (-2N) = 0$.

Verified to machine precision ($|\delta\lambda_m| < 10^{-14}$)
for $4 \le N \le 12$
(\texttt{proofs/rg\_stability\_characterization.py}).
\end{proof}

\begin{remark}[RG classification versus stability classification]
\label{rem:rg-vs-stab}
Proposition~\ref{prop:alpha-2-marginal} shows that the
\emph{stability-marginal} exponent ($\alpha = 2$) differs from
the \emph{RG-marginal} exponent ($\alpha = 0$).  The RG
classification (Proposition~\ref{prop:rg-classification}) concerns
the behaviour of the interaction under spatial coarse-graining;
the stability classification concerns the perturbation
$\delta\lambda_m$ of the Havelock eigenvalue.  Their relationship is:

\smallskip
\begin{center}
\begin{tabular}{ccccl}
\toprule
$\alpha$ & RG class & $A = \alpha(\alpha{-}2)$ &
  $\delta\lambda_3(N{=}7)$ & Stability effect \\
\midrule
$-2$ & irrelevant & $8$ & $-4\varepsilon$ & destabilising \\
$-1$ & irrelevant & $3$ & $-0.96\varepsilon$ & destabilising \\
$1$  & relevant   & $-1$ & $-0.58\varepsilon$ & destabilising \\
$2$  & relevant   & $0$ & $0$ & isotropic-marginal \\
$3$  & relevant   & $3$ & $+10.6\varepsilon$ & stabilising \\
$4$  & relevant   & $8$ & $+56\varepsilon$ & stabilising \\
\bottomrule
\end{tabular}
\end{center}
\smallskip

\noindent
All entries with $\alpha < 2$ have $\delta\lambda_3 < 0$
(destabilising for $\varepsilon > 0$); all entries with
$\alpha > 2$ have $\delta\lambda_3 > 0$ (stabilising for
$\varepsilon > 0$).  The crossover is at $\alpha = 2$, where
the anisotropy $A$ vanishes.
The Cauchy blob
($g = -r^{-2}$, i.e.\ $\alpha = -2$ with $\varepsilon < 0$)
stabilises the $N = 7$ marginal ring because the sign reversal
$\varepsilon = -\varepsilon_0^2$ converts $\delta\lambda_3
= (-4)(-\varepsilon_0^2) = 4\varepsilon_0^2 > 0$.
\end{remark}

\noindent\emph{Validity of first-order theory.}
The perturbation $\delta\lambda_m$ is first-order in $\varepsilon$
and reliable when $|\varepsilon\,\delta\lambda_m| \ll |\lambda_m(0)|$.
For the $N = 7$ marginal ring ($\lambda_3(0) = 0$), any nonzero
$\delta\lambda_3$ dominates regardless of $\varepsilon$, so the
\emph{sign} of $\delta\lambda_3$ is the reliable prediction.
For the $N = 8$ ring with $\alpha = 3$, the mode-resolved
computation gives $\delta\lambda_3 = 11.8$ and
$\delta\lambda_4 = 11.3$ (the two binding modes), so
$\varepsilon = 0.1$ predicts $\lambda_3 = -0.5 + 1.18 = 0.68 > 0$
and $\lambda_4 = -1.0 + 1.13 = 0.13 > 0$---all modes of the $8$-gon
become stable, giving $N_{\mathrm{crit}} = 8$ at first order.
Quantitative predictions for $N \ge 9$ require the full nonlinear
eigenvalue computation, not the linearised theory.

The Gaussian curvature $K$ is a destabilizing perturbation when
$K>0$ and a stabilizing one when $K<0$.  On the sphere ($K>0$) the
deviation from the logarithmic Green's function lowers
$N_{\mathrm{crit}}$ below~$7$; on the hyperbolic plane ($K<0$) it
adds a positive $O(|K|)$ correction $12|K|$ to the formerly-marginal
$\lambda_3$ eigenvalue~\eqref{eq:lam3_H2}.  The flat-plane case
$K=0$ is therefore the \emph{least stable} point of the
constant-curvature family with respect to polygon formation:

\subsection{M\"obius covariance of the logarithmic energy}
\label{sec:mobius-cov}

\begin{proposition}[M\"obius covariance]\label{prop:mobius}
Under a conformal map $f\colon\mathbb{C}\to\mathbb{C}$, the
logarithmic energy transforms as
\[
  H(f(z_1),\dots,f(z_N))
  = H(z_1,\dots,z_N) - \sum_{k=1}^{N}\tfrac{N-1}{2}\ln|f'(z_k)|.
\]
Consequently, the point vortex equations of motion are
M\"obius-invariant: the single-site correction
$\tfrac{N-1}{2}\ln|f'(z_k)|$ depends only on~$z_k$, so
$\partial/\partial \bar z_j$ for $j \ne k$ annihilates it, and the
Hamiltonian gradient $\partial H/\partial \bar z_j$ is unchanged.
\end{proposition}
\begin{proof}
The identity $|f(z){-}f(w)| = |f'(z)|^{1/2}|f'(w)|^{1/2}|z{-}w|
\cdot(1 + O(|z{-}w|))$ is exact when $f$ is M\"obius (the error
term vanishes), so $\ln|f(z_j){-}f(z_k)| = \ln|z_j{-}z_k|
+ \tfrac12\ln|f'(z_j)| + \tfrac12\ln|f'(z_k)|$.  Summing over
all $\binom{N}{2}$ pairs, each $\ln|f'(z_k)|$ appears $N{-}1$
times.  Since $\tfrac12\ln|f'(z_k)|$ depends on $z_k$ alone,
$\partial[\sum_k\tfrac{N-1}{2}\ln|f'(z_k)|]/\partial\bar z_j
= \tfrac{N-1}{2}\,\partial\ln|f'(z_j)|/\partial\bar z_j$,
which is the single-site term in the $j$-th equation of motion.
For M\"obius $f$, $f'$ is rational and $\ln|f'|$ is harmonic away
from poles, so this term contributes a gauge-like shift to the
Hamiltonian that does not couple distinct vortices.
\end{proof}

\begin{corollary}[Dilation subgroup = scale-invariant Hamiltonian]
\label{cor:dilation-rg}
Among all M\"obius transformations, the dilation subgroup
$\{f(z) = bz : b > 0\}$ is uniquely characterised by having
constant derivative: $f'(z) = b$ for all~$z$.  Substituting
into Proposition~\ref{prop:mobius} gives a correction that is
\emph{position-independent}:
\begin{equation}\label{eq:dilation-rg}
  H(bz_1,\dots,bz_N)
  = H(z_1,\dots,z_N) - \tfrac{N(N-1)}{2}\ln b.
\end{equation}
This is the $N$-body form of the dilation-invariance condition of
Remark~\ref{rem:rg-fixed-point}: scaling changes $H$ only by an
additive constant, so the equations of motion
$\dot{\bar z}_j = \partial H/\partial z_j$ are dilation-invariant.
For any other M\"obius map $f$, the derivative $|f'(z_k)|$ varies
between sites, producing a configuration-dependent correction that
is not a pure constant; the Hamiltonian is \emph{not}
scale-invariant under general conformal maps, only under the
one-parameter dilation subgroup.

Equation~\eqref{eq:dilation-rg} is the $N$-body lift of the per-pair
uniqueness result in Remark~\ref{rem:rg-fixed-point}: the logarithmic
energy is the unique $N$-body Hamiltonian that is dilation-invariant.
\end{corollary}

\begin{remark}[The loxodromic thread]
\label{rem:loxodromic}
The logarithmic interaction and the logarithmic spiral are dual
manifestations of the same structure: the loxodromic one-parameter
subgroup $z \mapsto e^{(\sigma + i\alpha)t}z$ of the M\"obius group
acting on~$\mathbb{C}$.
\begin{itemize}
  \item The \emph{interaction} is logarithmic because it is the
    unique dilation-invariant pairwise potential
    (Remark~\ref{rem:rg-fixed-point}; Cauchy's functional equation).
  \item The natural \emph{deformation curves} are logarithmic spirals
    because they are orbits of the dilation-rotation generator
    $\sigma + i\alpha$; the centered spiral deformation (Paper~II, eq.~7.1) is precisely
    this orbit applied
    to the $N$-gon.
  \item The \emph{stability threshold} $N = 7$ is scale-independent
    because the Havelock eigenvalue $\lambda_m = (N{-}1) - m(N{-}m)/2$
    has no $R$-dependence---a direct consequence of
    equation~\eqref{eq:dilation-rg}.
  \item The \emph{palindromic} structure of the hyperbolic threshold
    polynomial $A\xi^2 + B\xi + A = 0$
    (equation~\eqref{eq:palindromic}) encodes the conformal inversion
    $\xi \leftrightarrow 1/\xi$ of the Poincar\'e disk, which is a
    M\"obius map exchanging the interior and exterior of the unit
    circle.
  \item The logarithm is the \emph{IR fixed point} of the
    coarse-graining flow on interactions
    (Proposition~\ref{prop:rg-classification}): compact vorticity
    distributions interact via $-\ln r + O(\varepsilon^2/r^2)$
    at large separation, with the correction irrelevant under
    rescaling.  The sign rule demarcates the basin of attraction.
\end{itemize}
These five consequences---the interaction, the deformation, the
threshold, the algebraic structure, and the RG flow---all descend
from M\"obius covariance (Proposition~\ref{prop:mobius}).
The unifying role of the M\"obius group is not merely a symmetry
of the equations of motion; it is the reason the logarithmic
interaction produces polygons at all.
\end{remark}

\begin{remark}[From sign rule to physical stability]
\label{rem:signrule-to-stability}
Theorem~\ref{thm:sign-rule} establishes $H''(0) < 0$ along the spiral
direction---an \emph{unconstrained} energy curvature.  This does not by
itself imply instability, because the spiral deformation changes the
conserved angular impulse $L = \sum|z_k|^2$ at $O(\sigma^2)$.  The
physical dynamics is constrained to $\{L = \mathrm{const}\}$, and
the Lagrange multiplier $\mu_L < 0$ shifts the curvature by
$-2\mu_L > 0$, converting the unconstrained saddle into a
\emph{constrained minimum} (the constrained energy minimum (Paper~II, Theorem~8.2) below).
The sign rule therefore plays an indirect but essential role: it
guarantees that the $N$-gon is a saddle of the unconstrained energy
with a specific sign structure ($H''(0) < 0$ along the spiral,
$H''(0) > 0$ along other directions), and the Lagrange shift converts
this saddle into a minimum for $N \le 7$.  For $N \ge 8$, the negative
curvature along the $m = \lfloor N/2\rfloor$ Fourier mode overwhelms
the Lagrange shift, and the constrained configuration is a saddle.
\end{remark}

%% ====================================================================

%% ====================================================================
\section{The sign rule for general pairwise interactions}
\label{sec:sign-rule}

\begin{theorem}[Sign rule for general interactions]\label{thm:sign-rule}
For $N \ge 3$ vortices on a ring with general pairwise interaction
$h(r)$, the energy along the centered spiral satisfies
\begin{equation}\label{eq:Hpp-general}
  H''(0) = \sum_{m=1}^{N-1}\Bigl[
    \bigl(d_m^2 h''(d_m) + d_m h'(d_m)\bigr) S_m
    + \frac{h'(d_m)\, m^2(N-m)}{N^2\, d_m}
  \Bigr],
\end{equation}
where $d_m = 2R\sin(\pi m/N)$ and
$S_m = \sum_{k=0}^{N-m-1}[(2k+m-(N{-}1))/(2N)]^2 \ge 0$.

If\, $h'(r) < 0$ and $(rh'(r))' \le 0$ for all $r > 0$, then
$H''(0) < 0$.
\end{theorem}

\begin{proof}
\emph{Step 1: Pairwise distance as a function of $\sigma$.}
For vortices $j$ and $k$ with separation $m = j - k > 0$, write
the centered exponent as $(j - (N{-}1)/2)/N = (j' + k')/2 + m/(2N)$
where $j' = (j-(N{-}1)/2)/N$ and $k' = (k-(N{-}1)/2)/N$.
Define the mean displacement $s = \sigma(j'+k')/2$ and half-separation
$t = \sigma m/(2N)$.  Then $z_j(\sigma) = R\,e^{s+t}\,e^{i\theta_j}$,
$z_k(\sigma) = R\,e^{s-t}\,e^{i\theta_k}$ with
$\theta_j - \theta_k = 2\pi m/N$, giving
\[
  D(\sigma) \;{:=}\; |z_j - z_k|^2
  = R^2 e^{2s}\bigl[e^{2t} + e^{-2t} - 2\cos(2\pi m/N)\bigr]
  = R^2 e^{2s}\bigl[2\cosh(2t) - 2\cos(2\pi m/N)\bigr].
\]
At $\sigma = 0$: $s = t = 0$, so
$D(0) = 2R^2[1 - \cos(2\pi m/N)] = 4R^2\sin^2(\pi m/N) = d_m^2$.

\emph{Step 2: First and second derivatives of $D$.}
Write $c_s = (j'+k')/2$ and $c_t = m/(2N)$, so
$s = \sigma c_s$, $t = \sigma c_t$.  Direct differentiation gives
\begin{align*}
  D'(0) &= 4c_s(1 - \cos\phi), \qquad \phi = 2\pi m/N, \\[3pt]
  D''(0) &= 8c_s^2(1 - \cos\phi) + 8c_t^2.
\end{align*}

\emph{Step 3: Chain rule for $h(\sqrt{D})$.}
Write $d = \sqrt{D}$.  Two applications of the chain rule give
\[
  \frac{\dd^2}{\dd\sigma^2}h(d)\bigg|_{\sigma=0}
  = \underbrace{\frac{h''(d_m)}{4d_m^2}(D')^2
  + \frac{h'(d_m)}{2d_m}\!\left[D'' - \frac{(D')^2}{2d_m^2}\right]}_{\text{evaluated at each pair $(j,k)$}}.
\]
Substituting $D' = 4c_s d_m^2/2$ (since $1-\cos\phi = d_m^2/(2R^2)$,
and absorbing $R$ into $d_m$) and collecting terms:
\[
  \frac{\dd^2}{\dd\sigma^2}h(d)\bigg|_0
  = \bigl[d_m^2 h''(d_m) + d_m h'(d_m)\bigr]c_s^2
  + \frac{h'(d_m)}{d_m}\,c_t^2\,\frac{d_m^2}{(1-\cos\phi)}.
\]

\emph{Step 4: Summation over pairs.}
For fixed separation $m$, sum over the $N-m$ pairs
$(k, k{+}m)$ for $k = 0, \ldots, N{-}m{-}1$.  Define
\[
  S_m = \sum_{k=0}^{N-m-1} c_{s,k}^2, \quad
  \text{where } c_{s,k} = \frac{2k + m - (N{-}1)}{2N}.
\]
Since $\sum_k c_{s,k} = 0$ by the centering, $S_m \ge 0$ with equality
only if all $c_{s,k} = 0$ (which requires $N - m = 1$, i.e.\ $m = N-1$).
The second sum contributes $c_t^2 \cdot (N{-}m) = m^2(N{-}m)/(4N^2)$,
and $d_m^2/(1{-}\cos\phi) = 4R^2$.  Summing over $m = 1, \ldots, N-1$
yields \eqref{eq:Hpp-general}.

\emph{Step 5: Sign determination.}
Under $h'(r) < 0$: the second term
$B_m = h'(d_m)\,m^2(N{-}m)/(N^2 d_m) < 0$
for every $m = 1, \ldots, N-1$, since $h' < 0$, $d_m > 0$,
and $m^2(N{-}m) > 0$.
Under $(rh'(r))' \le 0$: the first coefficient
$A_m = d_m^2 h''(d_m) + d_m h'(d_m) = d_m \cdot (rh')'\big|_{r=d_m} \le 0$,
so the first term $A_m S_m \le 0$ since $S_m \ge 0$.

The case $m = N{-}1$ requires separate attention.
At $m = N{-}1$ there is only one pair $(0, N{-}1)$, and the
centering sum has a single term:
$c_{s,0} = (0 + (N{-}1) - (N{-}1))/(2N) = 0$,
so $S_{N-1} = 0$.  The first term $A_{N-1} S_{N-1} = 0$ therefore
\emph{vanishes identically}, regardless of the sign of $A_{N-1}$.
However, the second term $B_{N-1} = h'(d_{N-1})\,(N{-}1)^2/(N^2 d_{N-1})$
remains strictly negative (since $h' < 0$).  Thus the $m = N{-}1$
contribution to $H''(0)$ is $B_{N-1} < 0$, and strict negativity
of $H''(0) = \sum_m (A_m S_m + B_m)$ follows from $A_m S_m \le 0$
for all $m$ and $B_m < 0$ for all $m$, with at least one term
strictly negative.
\end{proof}

\begin{remark}[Borderline status of the logarithmic interaction]
\label{rem:borderline}
The sufficient conditions of Theorem~\ref{thm:sign-rule} partition
interactions into three classes:
\begin{enumerate}
  \item[(i)] $h' < 0$ and $(rh')' < 0$: both terms in
    \eqref{eq:Hpp-general} are strictly negative.
    Examples: $h = -r^\alpha$ ($\alpha > 0$), for which
    $(rh')' = -\alpha^2 r^{\alpha-1} < 0$.
  \item[(ii)] $h' < 0$ and $(rh')' = 0$: the first term vanishes
    and only the second contributes.  For smooth $h$ on $(0,\infty)$,
    this class contains a unique functional form, up to positive
    scaling and additive constant:
    $(rh')' = 0 \Rightarrow rh' = c \Rightarrow h' = c/r
    \Rightarrow h = c\ln r + b$; combined with $h' < 0$ gives
    $c < 0$, i.e.\ $h(r) = -a\ln r + b$ with $a > 0$.
  \item[(iii)] $h' < 0$ and $(rh')' > 0$ at some $r$: the sign of
    the first term is indefinite and $H''(0) < 0$ is not guaranteed.
\end{enumerate}
The 2D Green's function $h(r) = -\ln r$ falls in class~(ii)---the
boundary between guaranteed and non-guaranteed negativity.  In this
precise sense, the logarithmic interaction produces the \emph{weakest}
possible guarantee of $H''(0) < 0$: only the $B_m$ terms contribute
(the $A_m$ terms vanish identically), whereas any class~(i)
interaction has both terms strictly negative.  The logarithm is the
least robust polygon producer among all interactions satisfying the
sign rule.
\end{remark}

%% ====================================================================
\section{The Havelock identity and its metric-independence}
\label{sec:havelock-identity}

\subsection{The classical identity}
\label{sec:classical-havelock}
\label{app:havelock}

\begin{lemma}\label{lem:havelock}
For integers $N \ge 2$ and $1 \le m \le N-1$,
\[
  T_m := \sum_{p=1}^{N-1}\frac{1-\cos(2\pi pm/N)}{4\sin^2(\pi p/N)}
  = \frac{m(N-m)}{2}.
\]
\end{lemma}

This identity appears in \citet{Havelock1931}, \S3, as a step in
computing the stability eigenvalues of the $N$-vortex ring; see also
\citet{Aref2003}, Appendix~A.

\begin{proof}
Write $\omega = e^{2\pi i/N}$.  Using $2\sin(\pi p/N) =
|1 - \omega^p|$, expand
$[1 - \cos(2\pi pm/N)] / [4\sin^2(\pi p/N)]
= \frac12\bigl|\sum_{k=0}^{m-1}\omega^{pk}\bigr|^2
= \frac{1}{2}\sum_{j=0}^{m-1}\sum_{k=0}^{m-1}\omega^{p(j-k)}$.
Summing over $p = 1, \ldots, N-1$ and using the orthogonality relation
$\sum_{p=1}^{N-1}\omega^{p\ell} = N\delta_{\ell \equiv 0\pmod{N}} - 1$:
\[
  T_m = \frac{1}{2}\sum_{j,k=0}^{m-1}
  \bigl(N\delta_{j \equiv k\pmod{N}} - 1\bigr)
  = \frac{1}{2}(Nm - m^2) = \frac{m(N-m)}{2},
\]
where $\delta_{j \equiv k}$ contributes $m$ (since $0 \le j,k \le m-1
< N$ forces $j = k$) and the $-1$ contributes $-m^2$.
\end{proof}

\subsection{The Riemannian Havelock identity}
\label{sec:riemannian}

\begin{theorem}[Metric-independence of $f(m,N)$]
\label{thm:riemannian-havelock}
Let $M$ be a two-dimensional Riemannian manifold (compact or
non-compact) on which the Laplace--Beltrami Green's function $G$
depends only on geodesic distance: $G(x,y) = h(d(x,y))$.
Let the $N$-gon at geodesic radius $R$ be an equilibrium of the
point vortex Hamiltonian $H = -\sum_{j<k}G(z_j,z_k)$.
Then the constrained Lagrangian eigenvalues satisfy
\begin{equation}\label{eq:riemannian-havelock}
  \lambda_m = C_1(M, R) - \frac{m(N-m)}{2},
  \qquad m = 1, \ldots, N-1,
\end{equation}
where $C_1$ depends on the surface and ring radius but
$f(m,N) = m(N{-}m)/2$ is universal.
\end{theorem}

\noindent
The hypothesis ``$G$ depends only on geodesic distance'' holds on:
\begin{itemize}
  \item the plane ($G = -(2\pi)^{-1}\ln r$, $C_1 = N{-}1$);
  \item the sphere $\mathbf{S}^2$ of radius $a$
    ($G = -(4\pi)^{-1}\ln(2{-}2\cos\gamma)$,
    $C_1 = (N{-}1)(1{-}\xi)/(1{+}\xi)$, $\xi = KR^2$);
  \item the hyperbolic plane $\mathbf{H}^2$ of curvature
    radius $a$
    ($G = -(2\pi)^{-1}\ln\tanh(d/2)$,
    $C_1 = (N{-}1)(1{+}\xi^2)/(1{-}\xi)^2$, $\xi = r_E^2/a^2$).
\end{itemize}
It does \emph{not} hold on surfaces of non-constant curvature
(where $G$ depends on the positions individually, not just on
their separation) or on generic quotients $\Gamma\backslash M$
(where $G$ is the automorphic Green's function, depending on all
images $\gamma(y)$).

\begin{proof}
The proof uses two ingredients, both independent of $G$.

\emph{Step 1: Fourier-block structure.}
Let $\mathbf{M}_p$ be the $2{\times}2$ Hessian of $h(d_{0p})$
in the local (radial, tangential) frame at vortex~$0$.
The $\mathbb{Z}_N$-Fourier decomposition of the full Hessian gives,
for each component $\alpha,\beta\in\{r,t\}$:
\[
  H_{\alpha\beta}^{(m)}
  = \sum_{p=1}^{N-1} M_{p,\alpha\beta}
    \bigl(\cos(2\pi pm/N)-1\bigr),
\]
since the diagonal (self-interaction) block equals $-\sum_p M_p$
(the equilibrium condition forces the total gradient to vanish).

\emph{Step 2: Mode-independent trace.}
Summing the radial and tangential components:
\begin{equation}\label{eq:trace-formula}
  H_{rr}^{(m)} + H_{tt}^{(m)}
  = \sum_{p=1}^{N-1} \nabla^2_{\mathrm{LB}} h(d_p)
    \bigl(\cos(2\pi pm/N) - 1\bigr),
\end{equation}
where $\nabla^2_{\mathrm{LB}} h(d_p) =
\mathrm{tr}[\mathbf{M}_p]$.  The key input:
away from the diagonal ($x \ne y$), the Laplace--Beltrami
equation gives $\nabla^2_{\mathrm{LB},x}\,G(x,y) = c$, where
\begin{equation}\label{eq:laplacian-constant}
  c = \begin{cases}
    0 & \text{plane or $\mathbf{H}^2$ (non-compact, $\mathrm{Vol}=\infty$)}, \\
    -1/\mathrm{Vol}(M) & \text{compact surface (from $\Delta G = \delta_y - 1/\mathrm{Vol}$)}.
  \end{cases}
\end{equation}
Since $d_p \ne 0$ for all pairs on the $N$-gon (no two vortices
coincide), $\nabla^2_{\mathrm{LB}} h(d_p) = c$ for every $p$.
Substituting into~\eqref{eq:trace-formula}:
\[
  H_{rr}^{(m)} + H_{tt}^{(m)}
  = c \sum_{p=1}^{N-1}(\cos(2\pi pm/N) - 1)
  = c(-1 - (N{-}1)) = -Nc
  \eqqcolon \delta,
\]
using $\sum_{p=1}^{N-1}\cos(2\pi pm/N) = -1$ for $m \ge 1$.
The trace $\delta$ is independent of~$m$: it enters $C_1$ as an
additive shift but leaves $f(m,N) = m(N{-}m)/2$ unchanged.

On the plane: $c = 0$, $\delta = 0$ (exact tracelessness).
On $\mathbf{S}^2$ of radius $a$:
$c = -1/(4\pi a^2)$, $\delta = N/(4\pi a^2)$.
On $\mathbf{H}^2$: $c = 0$ (non-compact), $\delta = 0$.
On a compact quotient $\Gamma\backslash\mathbf{H}^2$:
$c = -1/\mathrm{Vol}(\Gamma\backslash\mathbf{H}^2)$,
$\delta = N/\mathrm{Vol}$---but note that on such a quotient
the Green's function is automorphic and does NOT depend only on
geodesic distance, so the hypothesis of the theorem is not
satisfied.  The identity~\eqref{eq:riemannian-havelock}
applies to the universal cover, not to the quotient.
\end{proof}

The stability threshold $N_{\mathrm{crit}}$ is therefore determined by
$C_1(\text{surface}) \ge \max_m m(N{-}m)/2$: a condition involving the
\emph{universal} function $f$ and a single surface-dependent constant.
On the plane: $C_1 = N{-}1$, giving $N_{\mathrm{crit}} = 7$ (the
largest integer with $(N{-}1) \ge \lfloor N/2\rfloor(N{-}\lfloor
N/2\rfloor)/2$).  On the sphere:
$C_1(\varphi_0) < N{-}1$ for $\varphi_0 > 0$, giving the
colatitude-dependent threshold of Section~\ref{sec:curved}.

\begin{remark}[Representation-theoretic content of the stability boundary]
\label{rem:rep-theory}
The decomposition $\lambda_m = C_1 - f(m,N)$ separates the stability
condition into a $\mathbb{Z}_N$ representation-theoretic part $f$ and
a geometric part $C_1$, with three consequences.

\emph{1.\ Spectral gap as a quantum phase transition.}
Write $H_N = C_1 I - F_N$ where $F_N = \mathrm{diag}(f(0,N),
\ldots, f(N{-}1,N))$ in the Fourier basis.  Stability is equivalent
to the positivity condition $H_N \ge 0$.  The spectral gap
$\Delta(N) = C_1 - f_{\max}(N)$ closes at $N = 7$ on the flat
plane ($\Delta = 0$), marking a transition from the gapped phase
($N \le 6$, all eigenvalues positive) to the broken phase
($N \ge 8$, negative eigenvalue).  Gaussian curvature acts as
a tuning parameter: $K < 0$ opens the gap, $K > 0$ closes it.

\emph{2.\ Crystallographic restriction analogue.}
The classical crystallographic restriction constrains the order
of a rotation compatible with a lattice: $\mathrm{tr}(R) = 2\cos(2\pi/N)
\in \mathbb{Z}$ forces $N \in \{1,2,3,4,6\}$ in~$\mathbb{R}^2$.
Our stability restriction constrains the order of a $\mathbb{Z}_N$
polygon compatible with the geometry: $C_1 \ge f_{\max}(N)$ forces
$N \le 7$ on $\mathbb{R}^2$.  Both are discrete constraints on
$N$ arising from the interplay between $\mathbb{Z}_N$ algebra
(trace / Casimir) and the ambient geometry (lattice / curvature).

\emph{3.\ Spectral statistics and Landau transition.}
Level spacings $(N{-}2m{-}1)/2$ are arithmetic (maximally rigid).
The gap closure $\Delta(N) \sim -(N{-}7)$ has mean-field exponents
($\nu = 1$, $\beta = 1/2$).  The quartic coefficient
$\alpha_0 = 45/14 > 0$ that resolves the marginal $N = 7$ case
is established in companion Paper~II, \S10; here we note only that
its positivity is predicted by the spectral structure (the gap
closes linearly, requiring a quartic restoring term).
\end{remark}

The value $C_1 = N{-}1$ on the plane has a geometric origin: for
$h' = -1/d$ (the logarithmic force law), each vortex pair $(0,k)$
contributes the \emph{same} amount $-1/(2R)$ to the gradient
$\partial H/\partial x_0$, regardless of the separation~$k$.  This
is because $h'(d_k)(x_0{-}x_k)/d_k = -(1{-}\cos(2\pi k/N))/
(4R\sin^2(\pi k/N)) = -1/(2R)$: the geometric factor
$(1{-}\cos)$ cancels against the denominator $\sin^2 = (1{-}\cos)/2$.
The cancellation is the \emph{scale invariance} of $h' = -1/d$:
the force law $h' \propto 1/d$ precisely compensates the
$d$-dependence of the chord geometry, making each pair's contribution
independent of the separation.  This property characterizes the
logarithm \emph{uniquely}: among radially symmetric pairwise
interactions on $\mathbb{R}^2$, $h = -\ln|z{-}w|$ is the unique
function satisfying $h(d\lambda) = h(d) + h(\lambda)$ (additivity
under rescaling---Cauchy's functional equation with continuity).
This additivity under the dilation subgroup $z \mapsto \lambda z$
(for which $|f(z){-}f(w)| = |\lambda|\,|z{-}w|$) is Cauchy's
multiplicative functional equation $h(ab) = h(a) + h(b)$,
whose unique continuous solution is $h = c\ln$.


%% ====================================================================
\section{Stability thresholds on constant-curvature surfaces}
\label{sec:curvature-thresholds}
\label{sec:curvature-invariant}

\noindent\emph{Prior work and what is new.}
\citet{BoattoCabral2003} computed the constrained Lagrangian eigenvalues
for a latitudinal ring of point vortices on the non-rotating sphere at
arbitrary colatitude, establishing that the stability boundary shifts
with latitude.  \citet{LimMontaldiRoberts2001} studied relative
equilibria of vortex rings on $S^2$ within the Hamiltonian framework.
The spherical Havelock eigenvalues are therefore not novel.

What the proof in~\S\ref{sec:riemannian} adds is an explicit algebraic
argument that the $m$-dependent part $f(m,N) = m(N{-}m)/2$ is
\emph{metric-independent}---a consequence of $\mathbb{Z}_N$
representation theory alone---while only the additive constant $C_1$
carries the geometric information.  A specialist who worked through the
Boatto--Cabral calculation would likely have noticed this; the value
here is in making the separation explicit and, more substantively, in
extending the framework to the hyperbolic plane~(\S\ref{sec:hyperbolic}),
where $N_{\mathrm{crit}} > 7$ and systematic stability tables do not
appear to be in the existing literature.

\subsection{Curved 2D surfaces: the spherical stability boundary}
\label{sec:curved}

On the 2-sphere, the point vortex interaction
$h(\theta) = -\ln(2\sin(\theta/2))$ satisfies $h' < 0$ and
$(\sin\theta\cdot h')' = \tfrac12\sin\theta > 0$ (class~(iii)).
While $H''(0) > 0$ along the spiral at all tested colatitudes
(the polygon is a spiral-direction energy minimum even without
the Lagrange shift), the \emph{constrained} Hessian tells a
different story.

The angular impulse on the sphere is
$L = \sum\cos\theta_k$, with Hessian
$\partial^2 L/\partial\theta_k^2 = -\cos\varphi_0$ at colatitude
$\varphi_0$.  The Lagrangian shift from $\mu_L$ is therefore
proportional to $\cos\varphi_0$: full near the pole ($\varphi_0 \to 0$),
vanishing at the equator ($\varphi_0 = 90^\circ$).  This
latitude-dependent shift produces a \emph{colatitude-dependent
stability boundary}:

\begin{center}
\begin{tabular}{lcccccc}
\toprule
$N$ & 3 & 4 & 5 & 6 & 7 & $\ge 8$ \\
\midrule
$\varphi_{\mathrm{crit}}$ & $60.0^\circ$ & $48.2^\circ$ &
  $41.4^\circ$ & $25.8^\circ$ & $0^\circ$ & --- \\
\bottomrule
\end{tabular}
\end{center}

\noindent
The $N$-gon is a constrained energy minimum on the sphere if and only
if $\varphi_0 < \varphi_{\mathrm{crit}}(N)$.  The dimensionless
eigenvalue $\bar\lambda = \sin^2\!\varphi_0 \cdot \lambda$ recovers
the Havelock values to within $1\%$ at $\varphi_0 = 3^\circ$
(N=5: $\bar\lambda = 0.99$, Havelock $= 1.0$;
N=6: $0.49$ vs $0.5$; N=8: $-1.01$ vs $-1.0$).

\emph{Illustration.}
Saturn's hexagon at $76^\circ$\,N ($\varphi_0 = 14^\circ$) lies
inside $\varphi_{\mathrm{crit}}(6) = 25.8^\circ$: the hexagonal
ring is in the stable zone (the planetary analysis is in Paper~III;
here we note only the mathematical consistency).
For $N = 7$: $\cos\varphi_{\mathrm{crit}} = 6/6 = 1$, so
$\varphi_{\mathrm{crit}} = 0^\circ$---the $7$-gon is marginally
stable only at the exact pole, consistent with the flat-plane
marginality $\lambda_3 = 0$.
For $N \ge 8$: $\cos\varphi_{\mathrm{crit}} > 1$, which is
unattainable; no stable latitude exists on the sphere.
Jupiter's octagon ($N = 8$) is therefore unstable at all
colatitudes, consistent with the need for a central vortex.

\paragraph{Derivation of $C_1(\mathbf{S}^2)$ and exact thresholds.}
For $N$ vortices on a latitude circle at colatitude $\varphi_0$, the
great-circle separation satisfies
$2\sin(\gamma_{jk}/2) = 2\sin\varphi_0\,|\sin(\Delta\phi_{jk}/2)|$,
so the spherical pair interaction decomposes as
\[
  -\ln\bigl(2\sin(\gamma_{jk}/2)\bigr)
  = -\ln\bigl(2|\sin(\Delta\phi_{jk}/2)|\bigr) - \ln(\sin\varphi_0).
\]
The first term is the flat Thomson energy in the azimuthal variable
$\phi$; its Hessian along azimuthal Fourier modes gives the Havelock
eigenvalues $(N{-}1)/2 - m(N{-}m)/2$ exactly as in the flat case.
The second term $-\binom{N}{2}\ln(\sin\varphi_0)$ is purely
radial and enters through the Lagrange multiplier.

The angular impulse on $\mathbf{S}^2$ is $L = \sum\cos\varphi_k$.
At the ring, $\partial^2 L/\partial\varphi_k^2 = -\cos\varphi_0$,
so the Lagrange shift
$\mu_L\cdot\partial^2 L/\partial\varphi_k^2$ is proportional to
$\cos\varphi_0$: maximal near the pole, vanishing at the equator.
Computing $\mu_L$ from the stationarity condition and combining with
the radial Hessian of the pair interaction yields the curvature
coefficient (cf.\ \citealt{BoattoCabral2003}):
\begin{equation}\label{eq:C1_S2}
  C_1(\mathbf{S}^2,\,\xi) = \frac{(N-1)(1-\xi)}{1+\xi} = (N-1)\cos\varphi_0,
\end{equation}
where $\xi = \tan^2(\varphi_0/2)$ is the stereographic parameter and
$\cos\varphi_0 = (1{-}\xi)/(1{+}\xi)$.
Marginal stability $C_1 = m(N{-}m)/2$ therefore reduces to
\begin{equation}\label{eq:cos_crit}
  \cos\varphi_{\mathrm{crit}} = \frac{m(N{-}m)}{2(N-1)}, \qquad m = \lfloor N/2\rfloor,
\end{equation}
which is \emph{always rational}.  The threshold in $\xi$ is
\begin{equation}\label{eq:xi_crit_S2}
  \xi_{\mathrm{crit}}(N) = \frac{(N-1) - m(N{-}m)/2}{(N-1) + m(N{-}m)/2}
  \in \mathbb{Q}.
\end{equation}
Explicit values:

\begin{center}
\begin{tabular}{lccccc}
\toprule
$N$ & $m(N{-}m)/2$ & $\cos\varphi_{\mathrm{crit}}$ & $\xi_{\mathrm{crit}}$ &
  $\varphi_{\mathrm{crit}}$ & stable iff \\
\midrule
3 & $1$ & $1/2$ & $1/3$ & $60.0^\circ$ & $\varphi_0 < 60^\circ$ \\
4 & $2$ & $2/3$ & $1/5$ & $48.2^\circ$ & $\varphi_0 < 48.2^\circ$ \\
5 & $3$ & $3/4$ & $1/7$ & $41.4^\circ$ & $\varphi_0 < 41.4^\circ$ \\
6 & $9/2$ & $9/10$ & $1/19$ & $25.8^\circ$ & $\varphi_0 < 25.8^\circ$ \\
7 & $6$ & $1$ & $0$ & $0^\circ$ &
  pole only (marginal) \\
$\ge 8$ & $> N{-}1$ & $> 1$ & --- & --- & never \\
\bottomrule
\end{tabular}
\end{center}

\noindent
The thresholds $\xi_{\mathrm{crit}} \in \{1/3,\,1/5,\,1/7,\,1/19\}$ are exact
rational numbers even though the Green's function $h(\theta) =
-\ln(2\sin(\theta/2))$ is transcendental---the rationality is a consequence of
the linear dependence of $C_1$ on $\cos\varphi_0$ (\eqref{eq:C1_S2}) combined
with the rational marginal condition~\eqref{eq:cos_crit}.

\begin{remark}[Pattern $\xi_{\mathrm{crit}} = 1/(2N{-}3)$ and its break at $N=6$]
For $N = 3, 4, 5$ one has $m(N{-}m)/2 = N{-}2$ exactly (since $m = \lfloor
N/2\rfloor = (N-1)/2$ or $N/2$ gives the same result), so the numerator
$(N-1) - (N-2) = 1$ and denominator $(N-1) + (N-2) = 2N-3$, yielding
$\xi_{\mathrm{crit}} = 1/(2N{-}3)$.  At $N = 6$ this pattern breaks: the
critical mode has $m(N{-}m)/2 = 9/2 \ne N{-}2 = 4$, because the hexagonal
mode carries a larger Havelock eigenvalue ($\lambda_{\min} = 1/2$ vs.\ $1$
for $N \le 5$).  Consequently $\xi_{\mathrm{crit}}(6) = 1/19$, not the
$1/(2 \cdot 6{-}3) = 1/9$ that the pattern would predict.

The notation $\bar\lambda = \sin^2\!\varphi_0 \cdot \lambda$ (line~938--941)
converts the dimensionful eigenvalue $\lambda$ to the dimensionless form
used in the numerical table above; the exact $\xi$-threshold~\eqref{eq:xi_crit_S2}
is independent of this rescaling.

The colatitudes in the summary table above are the exact values
$\varphi_{\mathrm{crit}} = \arccos(m(N{-}m)/(2(N{-}1)))$;
for $N = 6$, $\cos\varphi_{\mathrm{crit}} = 9/10$.
\end{remark}

\noindent
All exact threshold computations are verified by the companion
software (Appendix~A, Code Availability).

On the flat torus, $H''(0) < 0$ approaches the planar value for
polygons small relative to the domain, with periodic corrections
entering at the domain scale.

\subsection{The hyperbolic plane and Gaussian curvature}
\label{sec:hyperbolic}

\citet{Kimura1999} and others have studied point vortex dynamics on
surfaces of constant curvature; the Hamiltonian for the hyperbolic plane
is known.  Systematic stability tables for $N$-gons on $\mathbf{H}^2$
do not appear to be in the literature, and the results of this subsection
are likely the most novel part of Section~\ref{sec:curvature-invariant}.

The threshold $N_{\mathrm{crit}}$ on any surface is the largest $N$
with $C_1(\text{surface}) \ge \max_m m(N{-}m)/2$.
Gaussian curvature $K$ is the control parameter:
\begin{center}
\begin{tabular}{lcl}
\toprule
Curvature & $N_{\mathrm{crit}}$ & Mechanism \\
\midrule
$K > 0$ (sphere) & $< 7$ & Positive curvature destabilizes \\
$K = 0$ (plane) & $= 7$ & Critical transition \\
$K < 0$ (hyperbolic) & $> 7$ & Negative curvature stabilizes \\
\bottomrule
\end{tabular}
\end{center}
\begin{remark}[Order of the curvature correction]
\label{rem:euler}
The curvature enters $N_{\mathrm{crit}}$ at $O(K)$, not $O(K^2)$:
it acts through the surface Green's function via the Havelock
coefficient $C_1(\text{surface},\xi)$
(equations~\eqref{eq:C1_S2},~\eqref{eq:C1_H2}), which is linear
in $K$ at leading order.  On a compact surface the Gauss--Bonnet
constraint $\int K\,\mathrm{d}A = 2\pi\chi$ enters only through
$C_1$; no additional $O(K^2)$ self-energy correction arises for
point vortices (such a correction would require finite-area patches,
where it is $O(\varepsilon^2)$ and topology-independent).
\end{remark}

\noindent
On the hyperbolic plane, $C_1(\xi) = (N{-}1)(1{+}\xi^2)/(1{-}\xi)^2$
(equation~\eqref{eq:C1_H2} below) is strictly increasing in
$\xi = r_E^2/a^2$.  Setting $C_1(\xi^*) = \max_m m(N{-}m)/2$ and
solving the palindromic quadratic
$(c{-}1)\xi^2 - 2c\,\xi + (c{-}1) = 0$
with $c = \max_m m(N{-}m)/(2(N{-}1))$
gives the \emph{exact} threshold for each transition:
\begin{equation}\label{eq:xi-star-exact}
  \xi^*(N) = \frac{c - \sqrt{2c{-}1}}{c - 1},
  \qquad c = \frac{m(N{-}m)}{2(N{-}1)},\;\; m = \lfloor N/2\rfloor.
\end{equation}
No floating-point computation is involved: $\xi^*$ is an
\emph{algebraic number} for every~$N$, and an algebraic
\emph{integer} exactly when $N \in \{8, 9, 11, 15, 23\}$
(Proposition~\ref{prop:alg-int}).
\begin{center}
\begin{tabular}{ccccl}
\toprule
$N$ & $c$ & $\xi^*$ (exact) & $R/a = \sqrt{\xi^*}$ & Field \\
\midrule
$8$  & $8/7$  & $(8{-}3\sqrt{7})/7 = 0.0627$ & $0.250$ &
  $\mathbb{Q}(\sqrt{7})$ \\
$9$  & $5/4$  & $5{-}2\sqrt{6} = 0.1010$ & $0.318$ &
  $\mathbb{Q}(\sqrt{6})$ \\
$10$ & $25/18$ & $1/7 = 0.1429$ & $0.378$ &
  $\mathbb{Q}$ (rational) \\
$11$ & $3/2$  & $3{-}2\sqrt{2} = 0.1716$ & $0.414$ &
  $\mathbb{Q}(\sqrt{2})$ \\
$12$ & $18/11$ & $0.2024$ & $0.450$ &
  $\mathbb{Q}(\sqrt{5/11})$ \\
$15$ & $2$    & $2{-}\sqrt{3} = 0.2679$ & $0.518$ &
  $\mathbb{Q}(\sqrt{3})$ \\
$19$ & $5/2$  & $1/3 = 0.3333$ & $0.577$ &
  $\mathbb{Q}$ (rational) \\
$23$ & $3$    & $(3{-}\sqrt{5})/2 = 0.3820$ & $0.618$ &
  $\mathbb{Q}(\sqrt{5})$ \\
\bottomrule
\end{tabular}
\end{center}
The Riemannian Havelock identity $\lambda_m = C_1(R/a) - m(N{-}m)/2$
is verified to hold (computed $C_1$ values are mode-independent to
numerical precision~$\sim 10^{-5}$).  The exact $\mathbf{H}^2$
Green's function yields
\begin{equation}
  C_1(\mathbf{H}^2,\,\xi)
    = \frac{(N-1)(1+\xi^2)}{(1-\xi)^2},
  \qquad \xi = \frac{r_E^2}{a^2} = |K|\,r_E^2 ,
  \label{eq:C1_H2}
\end{equation}
where $r_E$ is the Euclidean ring radius in the Poincar\'e disk and
$K = -1/a^2$.

\begin{proof}[Derivation of~\eqref{eq:C1_H2}]
Work in the Poincar\'e disk of radius $a$ ($K = -1/a^2$).
The point vortex Hamiltonian and angular impulse are
\citep{Kimura1999}
\begin{align}
  H &= -\!\!\sum_{j<k}\ln|z_j{-}z_k|
  + \tfrac{N{-}1}{2}\!\sum_k\ln(a^2{-}|z_k|^2), \label{eq:H-H2}\\
  J &= \sum_k \frac{a^2{+}|z_k|^2}{a^2{-}|z_k|^2}
  \qquad\text{(SO(2,1) moment map)}.
  \label{eq:J-H2}
\end{align}
The $N$-gon equilibrium has $z_k = r_E\,e^{2\pi i k/N}$ with
$\xi = r_E^2/a^2$.

\emph{Step 1: equilibrium rotation frequency $\Omega$.}
At vortex $0 = (r_E, 0)$, the radial gradients are
(using $\sum_{p=1}^{N-1}(x_0{-}x_p)/|z_0{-}z_p|^2 = (N{-}1)/(2r_E)$
from the flat-plane Havelock sum):
\begin{align*}
  \frac{\partial H}{\partial x_0}
  &= -\frac{N{-}1}{2r_E}
  - \frac{(N{-}1)\,r_E}{a^2 - r_E^2}, \\
  \frac{\partial J}{\partial x_0}
  &= \frac{4a^2\,r_E}{(a^2 - r_E^2)^2}.
\end{align*}
The equilibrium condition
$\partial H/\partial x_0 = \Omega\,\partial J/\partial x_0$
gives
\[
  \Omega = \frac{-(N{-}1)}{2r_E}
  \cdot\frac{(a^2{-}r_E^2)^2}{4a^2 r_E}
  \cdot\Bigl(1 + \frac{2r_E^2}{a^2{-}r_E^2}\Bigr)
  = \frac{-(N{-}1)(a^2{+}r_E^2)}{4a^2 r_E^2}.
\]

\emph{Step 2: Hessian Rayleigh quotient.}
By Theorem~\ref{thm:riemannian-havelock}, the constrained
eigenvalue is
$\lambda_m = -\Omega\,(\partial^2 J/\partial r^2)^{-1}
- v_m^\top\nabla^2\!H\,v_m$.
But since the trace $\delta = 0$ on $\mathbf{H}^2$ (non-compact,
$c = 0$ in~\eqref{eq:laplacian-constant}), the Fourier-block
trace is zero and
$v_m^\top\nabla^2\!H\,v_m = (m(N{-}m)/2 - C_1^H)/ r_E^2$
where $C_1^H$ absorbs all surface-dependent terms.

\emph{Step 3: assembling $C_1$.}
Writing $\Omega = -(N{-}1)(1{+}\xi)/(4a^2\xi)$ (with
$\xi = r_E^2/a^2$), the Lagrange multiplier contributes
$-\Omega \cdot r_E^2 \cdot (\partial^2 J/\partial r^2)$
at the equilibrium.  Computing
$\partial^2 J/\partial r^2 = 4a^2(a^2{+}3r_E^2)/(a^2{-}r_E^2)^3$
and combining the flat Havelock term $(N{-}1)$, the conformal
correction $-(N{-}1)\,2\xi/(1{-}\xi)^2$ from the
$\ln(a^2{-}|z|^2)$ term in~\eqref{eq:H-H2}, and the
moment-map contribution, one obtains after simplification
(each term is a rational function of $\xi$; the algebra
is straightforward but lengthy):
\[
  C_1 = (N{-}1)\,\frac{1+\xi^2}{(1-\xi)^2}.
\]
Verified by numerical Hessian to relative precision $10^{-5}$
for $3 \le N \le 12$ and $0 < \xi < 0.9$
(\texttt{extensions/h2\_stability.py}).
\end{proof}

Since $(1+\xi^2)/(1-\xi)^2 > 1$ for all $\xi > 0$,
one has $C_1(\mathbf{H}^2) > N{-}1$ for all $R/a > 0$, confirming the
stabilizing effect of negative curvature.

$N_{\mathrm{crit}}$ increases through integer jumps as $R/a$ grows,
with the $7 \to 8$ transition at $R/a \approx 0.52$ (more precisely,
the $m = \lfloor N/2\rfloor$ eigenvalue of the $N = 8$ ring crosses
zero near this value).
We \emph{conjecture} that $N_{\mathrm{crit}}(R/a)$ is non-decreasing.
Numerical evaluation of the constrained Hessian eigenvalues
confirms this for $R/a \le 2$.
A plausible mechanism is the transition from class~(ii) at
$R \ll a$ to class~(i) at $R \gg a$, but the Lagrange multiplier
also depends on $R/a$ in a non-trivial way, so the monotonicity
is not proven.

\paragraph{Near-$K=0$ behavior.}
Expanding~\eqref{eq:C1_H2} for $N=7$, $m=3$ (see
companion Appendix~A for the algebraic simplification), the
$m=3$ eigenvalue satisfies
\begin{equation}
  \lambda_3(N=7,\,\xi)
    = \frac{12\xi}{r_E^2(1-\xi)^2}
    = \frac{12|K|}{(1-\xi)^2}
    = 12|K| + O(K^2 r_E^2) > 0 .
  \label{eq:lam3_H2}
\end{equation}
Negative curvature stabilizes the $N=7$ ring at \emph{first} order in
$|K|$: $N_\mathrm{crit}(K)$ stays at $7$ for all $K \le 0$ in the
sense that the formerly-neutral mode becomes strictly stable.

The $7\to 8$ transition is governed by the $N=8$ ring reaching marginal
stability.  Setting $C_1(\mathbf{H}^2,\xi) = m(N{-}m)/2\big|_{m=4} = 8$
gives $7(1+\xi^2)/(1-\xi)^2 = 8$, which simplifies to
\begin{equation}
  \xi^2 - 16\xi + 1 = 0
  \;\Longrightarrow\;
  \xi^* = 8 - 3\sqrt{7} \approx 0.0627 ,
  \label{eq:xi_star}
\end{equation}
exact.  In terms of geodesic ring radius $\rho$, this corresponds to
$\rho/a \approx 0.52$ (or $|K|\rho^2 \approx 0.27$).
The coefficient
\begin{equation}
  \gamma \;\equiv\; \frac{1}{(\xi^*)^2} = 127 + 48\sqrt{7} \approx 254
  \label{eq:gamma}
\end{equation}
carries an algebraic interpretation.
The ring of integers of $\mathbb{Q}(\sqrt{7})$ has fundamental unit
$\varepsilon = 8 + 3\sqrt{7}$, since $(8+3\sqrt{7})(8-3\sqrt{7}) = 64-63 = 1$.
Therefore $\xi^* = 8-3\sqrt{7} = \varepsilon^{-1}$ is the
\emph{inverse of the fundamental unit}---the unique unit below~$1$
in $\mathbb{Z}[\sqrt{7}]$---and $\gamma = (\xi^*)^{-2} = \varepsilon^2$.
The verification identity $(8-3\sqrt{7})^2(127+48\sqrt{7})=1$ is not a
computational coincidence: it is the norm-$1$ property
$\varepsilon^{-2}\cdot\varepsilon^2=1$ in $\mathbb{Z}[\sqrt{7}]$.

The field $\mathbb{Q}(\sqrt{7})$ arises through a specific algebraic
chain: for the $N = 8$ transition, $A = -1$ and $B = 16$, giving
discriminant $\Delta = B^2 - 4A^2 = 256 - 4 = 252 = 2^2 \cdot 3^2
\cdot 7$, with squarefree part~$7$.  The~$7$ enters as the factor
$(N{-}1) = 7$ in the general formula $\Delta = (N{-}1)(N{-}2)^2$
(Proof of~\eqref{eq:field-pattern})---the flat-plane stability
boundary $N = 7$ leaves its algebraic shadow in the discriminant.
The substantive content is the
\emph{palindromic structure}: the threshold polynomial has the form
$A\,\xi^2 + B\,\xi + A = 0$ (same leading and trailing coefficient),
which forces $\xi_1\,\xi_2 = A/A = 1$---the two roots are
multiplicative inverses, so $\xi^*$ is automatically an algebraic
\emph{unit} (an element of norm~$1$ in $\mathbb{Z}[\sqrt{7}]$).
This palindromic property is not generic: it reflects the inversion
symmetry $\xi\leftrightarrow 1/\xi$ of the Poincar\'e disk model,
which in turn is the restriction to the real axis of the conformal
automorphism $z \mapsto a^2/\bar z$ of the disk.  The stability
threshold is a unit \emph{because} the hyperbolic energy is
conformally covariant (Proposition~\ref{prop:mobius}): the palindromic
polynomial encodes the same M\"obius symmetry at the level of
algebraic number theory.

%% ====================================================================
\section{Algebraic structure of the stability boundary}
\label{sec:algebraic-structure}

The $7\to 8$ transition is not an isolated result: the $N \to (N{+}1)$
transition threshold on $\mathbf{H}^2$ is always the smaller positive root
of a palindromic quadratic, derived as follows.
The stability condition $\lambda_m(\xi^*) = 0$
(equation~\eqref{eq:riemannian-havelock}) with
$C_1 = (N{-}1)(1{+}\xi^2)/(1{-}\xi)^2$
(equation~\eqref{eq:C1_H2}) gives
\[
  (N{-}1)\,\frac{1+\xi^2}{(1-\xi)^2} = \frac{m(N{-}m)}{2}.
\]
Multiplying both sides by $(1{-}\xi)^2$ and rearranging:
\[
  (N{-}1)(1{+}\xi^2) = \tfrac{m(N{-}m)}{2}(1{-}\xi)^2
  = \tfrac{m(N{-}m)}{2} - m(N{-}m)\xi + \tfrac{m(N{-}m)}{2}\xi^2.
\]
Collecting powers of $\xi$:
\begin{equation}\label{eq:palindromic}
  A\,\xi^2 + B\,\xi + A = 0,
  \quad
  A = (N{-}1) - \tfrac{m(N{-}m)}{2},
  \quad
  B = m(N{-}m),
  \quad
  m = \lfloor N/2\rfloor.
\end{equation}
The coefficient of $\xi^2$ equals the constant term (both are~$A$):
this palindromic structure arises because $C_1(\xi)$ has the
form $(1{+}\xi^2)/(1{-}\xi)^2$, whose numerator and denominator
are both even functions of~$\xi$ (invariant under
$\xi \to -\xi$), producing equal leading and trailing coefficients
when cleared of denominators.
The palindromic property forces $\xi_1\xi_2 = A/A = 1$ (the
product of the two roots), so $\xi^* = 1/\xi^{*\prime}$---the
conformal inversion identity.
Since $A < 0$ for $N \ge 8$ and $B > 0$, the smaller positive root is
$\xi^* = (B/2 - \sqrt{(B/2)^2 - A^2})\,/\,|A|$, and its field extension
is determined by the squarefree part of the discriminant
$(B/2)^2 - A^2 \in \mathbb{Z}$:

\begin{center}
\begin{tabular}{cclcc}
\toprule
$N$ & $\xi^*$ (approx.) & Field & $D$ & $\xi^*$ (exact form) \\
\midrule
7  & 0 & $\mathbb{Q}$ & 1 & $0$ (marginal) \\
8  & 0.0627 & $\mathbb{Q}(\sqrt{7})$ & 7 & $8 - 3\sqrt{7}$ \\
9  & 0.1010 & $\mathbb{Q}(\sqrt{6})$ & 6 & $5 - 2\sqrt{6}$ \\
10 & 0.1429 & $\mathbb{Q}$ & 1 & $1/7$ \\
11 & 0.1716 & $\mathbb{Q}(\sqrt{2})$ & 2 & \\
12 & 0.2024 & $\mathbb{Q}(\sqrt{11})$ & 11 & \\
13 & 0.2162 & $\mathbb{Q}(\sqrt{10})$ & 10 & \\
14 & 0.2355 & $\mathbb{Q}(\sqrt{13})$ & 13 & \\
15 & 0.2533 & $\mathbb{Q}(\sqrt{3})$ & 3 & \\
16 & 0.2699 & $\mathbb{Q}(\sqrt{15})$ & 15 & \\
\bottomrule
\end{tabular}
\end{center}

\noindent
The discriminant's squarefree part $D$ follows the pattern
\begin{equation}\label{eq:field-pattern}
  D = \begin{cases}
    \mathrm{sq\_free}(N-1) & N \text{ even}, \\
    \mathrm{sq\_free}(N-3) & N \text{ odd},
  \end{cases}
\end{equation}
for all $N \ge 8$.  The deeper content is that the palindromic
structure encodes the conformal inversion
$\xi \leftrightarrow 1/\xi$ of the Poincar\'e disk (as noted above);
the field pattern itself follows from elementary factorization.

\begin{proof}[Proof of~\eqref{eq:field-pattern}]
The palindromic quadratic $A\xi^2 + B\xi + A = 0$ has discriminant
$\Delta = B^2 - 4A^2$.  Using $A = (N{-}1) - m(N{-}m)/2$ and
$B = m(N{-}m)$:
\[
  \Delta = m^2(N{-}m)^2 - 4\bigl[(N{-}1) - m(N{-}m)/2\bigr]^2
  = 4(N{-}1)\bigl[m(N{-}m) - (N{-}1)\bigr].
\]

\emph{Even $N$ ($m = N/2$):}
$m(N{-}m) = N^2/4$, so
$m(N{-}m) - (N{-}1) = N^2/4 - N + 1 = (N{-}2)^2/4$.
Therefore
$\Delta = 4(N{-}1)(N{-}2)^2/4 = (N{-}1)(N{-}2)^2$.
The squarefree part of $\Delta$ equals
$\mathrm{sq\_free}(N{-}1)$, since $(N{-}2)^2$ is a perfect square.
The threshold $\xi^*$ involves $\sqrt\Delta$, so
$\xi^* \in \mathbb{Q}(\sqrt{\mathrm{sq\_free}(N{-}1)})$.

\emph{Odd $N$ ($m = (N{-}1)/2$):}
$m(N{-}m) = (N{-}1)(N{+}1)/4 = (N^2{-}1)/4$, so
$m(N{-}m) - (N{-}1) = (N^2{-}1)/4 - (N{-}1) = (N{-}1)(N{-}3)/4$.
Therefore
$\Delta = 4(N{-}1) \cdot (N{-}1)(N{-}3)/4 = (N{-}1)^2(N{-}3)$.
The squarefree part of $\Delta$ equals
$\mathrm{sq\_free}(N{-}3)$, since $(N{-}1)^2$ is a perfect square.
Hence $\xi^* \in \mathbb{Q}(\sqrt{\mathrm{sq\_free}(N{-}3)})$.
\end{proof}

\paragraph{Algebraic integers and the Pell equation.}
The threshold $\xi^*(N)$ is an algebraic integer---i.e., it
satisfies a \emph{monic} polynomial $\xi^2 - c\,\xi + 1 = 0$ with
$c \in \mathbb{Z}$---if and only if $B/A \in \mathbb{Z}$.

\begin{proposition}[Algebraic-integer classification]
\label{prop:alg-int}
The threshold $\xi^*(N)$ is an algebraic integer if and only if
$N \in \{8, 9, 11, 15, 23\}$.
\end{proposition}

\begin{proof}
For odd $N \ge 9$: $A = -(N{-}1)(N{-}7)/8$,
$B = (N{-}1)(N{+}1)/4$, so $B/A = -2(N{+}1)/(N{-}7)$.
This is an integer iff $(N{-}7) \mid 2(N{+}1)$.
Since $2(N{+}1) = 2(N{-}7) + 16$, the condition reduces to
$(N{-}7) \mid 16$.  The positive divisors of~$16$ giving odd
$N \ge 9$ are $N{-}7 \in \{2, 4, 8, 16\}$, i.e.,
$N \in \{9, 11, 15, 23\}$.
For even $N \ge 8$: only $N = 8$ satisfies $B/A \in \mathbb{Z}$
(verified directly: $B/A = -16$; for $N \ge 10$ the denominator
$|A| = (N^2{-}8N{+}8)/8$ grows faster than $B = N^2/4$ allows
integer divisibility).
\end{proof}

The differences between consecutive algebraic-integer values
are $\{1, 2, 4, 8\}$---powers of~$2$---reflecting the
binary structure of the divisors of~$16$.

For these five values, $\xi^*$ satisfies the Pell equation.
Since the monic polynomial $\xi^2 - c\,\xi + 1 = 0$ forces
$\xi\cdot\xi' = 1$, writing $\xi^* = a + b\sqrt{D}$ with
$a, b \in \mathbb{Z}$ (or half-integers for $D \equiv 1 \pmod{4}$)
gives $a^2 - D b^2 = 1$: the classical Pell equation.
Each threshold is an explicit power of the inverse fundamental
unit $\varepsilon$ of $\mathbb{Q}(\sqrt{D})$:
\begin{center}
\begin{tabular}{ccllcc}
\toprule
$N$ & $D$ & Fundamental unit $\varepsilon$ &
$\xi^*(N)$ & $c = -B/A$ & Norm$(\varepsilon)$ \\
\midrule
$8$  & $7$ & $8 + 3\sqrt{7}$       & $\varepsilon^{-1}$
  & $16$ & $+1$ \\
$9$  & $6$ & $5 + 2\sqrt{6}$       & $\varepsilon^{-1}$
  & $10$ & $+1$ \\
$11$ & $2$ & $1 + \sqrt{2}$\,(silver ratio) & $\varepsilon^{-2}$
  & $6$  & $-1$ \\
$15$ & $3$ & $2 + \sqrt{3}$        & $\varepsilon^{-1}$
  & $4$  & $+1$ \\
$23$ & $5$ & $\varphi = (1{+}\sqrt{5})/2$\,(golden ratio)
  & $\varphi^{-2}$ & $3$  & $-1$ \\
\bottomrule
\end{tabular}
\end{center}
The exponent is $k = 1$ when $\mathrm{Norm}(\varepsilon) = +1$
(the fundamental unit solves $x^2 - Dy^2 = 1$) and $k = 2$ when
$\mathrm{Norm}(\varepsilon) = -1$ (the fundamental unit solves
$x^2 - Dy^2 = -1$; only its square has norm~$+1$, as required
by the palindromic constraint $\xi\cdot\xi' = 1$).

\paragraph{The golden ratio as the terminal case.}
The last algebraic-integer threshold ($N = 23$) has
$\xi^*(23) = \varphi^{-2} = (3 - \sqrt{5})/2$, where
$\varphi = (1{+}\sqrt{5})/2$ is the golden ratio.  Its monic
polynomial $\xi^2 - 3\xi + 1 = 0$ has the smallest integer
trace ($c = 3$) compatible with two distinct positive roots;
$c = 2$ gives the degenerate $(\xi{-}1)^2 = 0$.  The golden
ratio is thus the algebraically \emph{simplest} non-trivial
palindromic unit.
Equivalently, $\varphi^{-1}$ is the unique algebraic unit
in $(0,1)$ satisfying $1 - x = x^2$, which gives
$\sum_{k=1}^{\infty}\varphi^{-k}/k
= -\ln(1{-}\varphi^{-1}) = -\ln\xi^*(23) = 2R(5)$:
the generating series of the Fibonacci weights evaluates to
twice the regulator.

\paragraph{Binary subdivision.}
The odd algebraic-integer values $N = 7 + 2^k$ ($k = 1, 2, 3, 4$)
have traces $c(k) = 2 + 2^{4-k}$: each doubling of $(N{-}7)$
halves the algebraic excess $c - 2$.  The sequence
$c \in \{10, 6, 4, 3\}$ converges to the golden-ratio value
$c = 3$ as $k \to 4$.  The binary subdivision terminates because
$c = 2$ is degenerate; no further algebraic-integer thresholds
exist beyond $N = 23$.

\paragraph{Continued fraction structure.}
\begin{proposition}[CF of palindromic units]\label{prop:cf-palindromic}
If\, $\xi$ satisfies $\xi^2 - c\,\xi + 1 = 0$ with integer
$c \ge 3$ and $\xi \in (0,1)$, then
$1/\xi = [c{-}1;\;\overline{1,\, c{-}2}]$.
For $c = 3$: the period degenerates to $[1]$ (all ones), which
is the continued fraction of the golden ratio~$\varphi$.
\end{proposition}

\begin{proof}
\emph{Step~1.}  $a_0 = \lfloor 1/\xi\rfloor = c - 1$, since
$(c{-}2)^2 \le c^2 - 4$ for $c \ge 2$.
\emph{Step~2.}  The first remainder is
$r_0 = (\sqrt{c^2{-}4} - (c{-}2))/2$; its reciprocal
$1/r_0 = (\sqrt{c^2{-}4} + (c{-}2))/(2(c{-}2))$ satisfies
$1 \le 1/r_0 < 2$ (from $4(2c{-}5)(c{-}2) > 0$ for $c \ge 3$),
so $a_1 = 1$.
\emph{Step~3.}  $1/r_1 = 1/\xi - 1$, giving $a_2 = c - 2$.
\emph{Step~4.}  $r_2 = 1/\xi - (c{-}1) = r_0$ (periodicity).
\end{proof}

Applied to the five algebraic-integer cases:
$1/\xi^*(8) = [15;\;\overline{1, 14}]$,
$1/\xi^*(9) = [9;\;\overline{1, 8}]$,
$1/\xi^*(11) = [5;\;\overline{1, 4}]$,
$1/\xi^*(15) = [3;\;\overline{1, 2}]$,
$1/\xi^*(23) = [2;\;\overline{1}] = [2; 1, 1, 1, \ldots]$
(the golden ratio).

The $N = 10$ case gives the unique rational threshold
$\xi^*(10) = 1/7 \in \mathbb{Q}$
($N{-}1 = 9 = 3^2$, so $D = 1$).  A second rational threshold
occurs at $N = 19$: $\xi^*(19) = 1/3$
($N{-}3 = 16 = 4^2$, so $D = 1$).  Every even transition lives
in at most a degree-$2$ extension of $\mathbb{Q}$; the tower
terminates at~$\mathbb{Q}$ whenever
$\mathrm{sq\_free}(N{-}1) = 1$ (even $N$) or
$\mathrm{sq\_free}(N{-}3) = 1$ (odd $N$).

\paragraph{The regulator as physical accessibility.}
The regulator $R(D) = \ln\varepsilon$ of $\mathbb{Q}(\sqrt{D})$
measures the ``size'' of the unit group.  For the five
algebraic-integer fields:
$R(7) = 2.77$, $R(6) = 2.29$, $R(3) = 1.32$, $R(2) = 0.88$,
$R(5) = 0.48$.
Since $\xi^* = \varepsilon^{-k}$, we have
$\ln(1/\xi^*) = k\,R(D)$: the regulator is the logarithmic
distance from the stability boundary to the flat-plane limit
$\xi = 0$.  Larger $R$ means smaller $\xi^*$, i.e., more
curvature is needed to reach the transition.
The $N = 8$ transition ($R = 2.77$, $\xi^* = 0.063$) requires
the most curvature; the $N = 23$ transition ($R = 0.48$,
$\xi^* = 0.382$) requires the least---making
$\varphi^{-2}$ the most experimentally accessible threshold.

\begin{proposition}[Galois group as conformal inversion group]
\label{prop:galois-inversion}
The composite field
$K = \mathbb{Q}(\sqrt{2}, \sqrt{3}, \sqrt{5}, \sqrt{6}, \sqrt{7})$
contains all five algebraic-integer thresholds.  Its Galois group
$\mathrm{Gal}(K/\mathbb{Q}) \cong (\mathbb{Z}/2)^4$ acts on the
threshold set by
$\sigma_D\colon \xi^*(N) \mapsto 1/\xi^*(N)$
---the conformal inversion of the Poincar\'e disk.
The Galois group of the stability threshold fields \emph{is} the
conformal inversion group, realised algebraically.
\end{proposition}

\begin{proof}
Each $\xi^*(N)$ lies in $\mathbb{Q}(\sqrt{D})$ for some squarefree
$D$ (Proposition~\ref{prop:alg-int}).  The Galois automorphism
$\sigma_D\colon \sqrt{D} \mapsto -\sqrt{D}$ sends $\xi^* = a + b\sqrt{D}$
to $a - b\sqrt{D} = 1/\xi^*$ (the palindromic conjugate).
Since the palindromic property encodes conformal inversion
$\xi \leftrightarrow 1/\xi$ (Section~\ref{sec:algebraic-structure}),
$\sigma_D$ acts as conformal inversion on each threshold.
The composite field $K$ has
$\mathrm{Gal}(K/\mathbb{Q}) \cong \prod_{D \in \{2,3,5,6,7\}}
\mathbb{Z}/2 \cong (\mathbb{Z}/2)^4$
(four independent generators, since $\sqrt{6} = \sqrt{2}\sqrt{3}$
is algebraically dependent), and each generator acts by conformal
inversion on the thresholds in its field component.
\end{proof}

\paragraph{Stability thresholds as closed geodesics.}
The monic palindromic polynomial $\xi^2 - c\,\xi + 1 = 0$
with $c \in \mathbb{Z}$ is simultaneously the characteristic polynomial
of a hyperbolic element $\gamma \in \mathrm{SL}(2,\mathbb{Z})$
with trace $\mathrm{tr}(\gamma) = c$.  The eigenvalues of $\gamma$
are $\varepsilon$ and $\varepsilon^{-1} = \xi^*$---exactly the
stability threshold and its conformal inverse.
Each such $\gamma$ defines a primitive conjugacy class of
$\mathrm{SL}(2,\mathbb{Z})$, hence a primitive closed geodesic
on the modular surface $\mathrm{SL}(2,\mathbb{Z})\backslash
\mathbf{H}^2$ of length
$\ell = 2\,\mathrm{arccosh}(c/2) = 2\ln\varepsilon$.

\begin{proposition}[Geodesic correspondence]
\label{prop:geodesic-correspondence}
The five algebraic-integer stability thresholds correspond to
five primitive closed geodesics on the modular surface.
For odd $N \in \{9, 11, 15, 23\}$ the traces are
$c = 2 + 16/(N{-}7)$
(equivalently, $N = 7 + 16/(c{-}2)$);
for $N = 8$ the trace is $c = 16$ (direct computation).
The full discriminant $\Delta = c^2 - 4$ decomposes as
$\Delta = f^2 D$ where $D$ is the squarefree part and $f$ is
the conductor:
\begin{center}
\begin{tabular}{ccccclll}
\toprule
$N$ & $c$ & $\Delta$ & $f$ & $D$ &
  Geodesic & $\mathrm{SL}(2,\mathbb{Z})$ \\
\midrule
$23$ & $3$  & $5$   & $1$ & $5$ &
  $\ell = 1.92$ & $\bigl[\begin{smallmatrix}2&1\\1&1\end{smallmatrix}\bigr]$
  \emph{(Fibonacci)} \\
$15$ & $4$  & $12$  & $2$ & $3$ &
  $\ell = 2.63$ & $\bigl[\begin{smallmatrix}3&2\\1&1\end{smallmatrix}\bigr]$ \\
$11$ & $6$  & $32$  & $4$ & $2$ &
  $\ell = 3.53$ & $\bigl[\begin{smallmatrix}5&2\\2&1\end{smallmatrix}\bigr]$ \\
$9$  & $10$ & $96$  & $4$ & $6$ &
  $\ell = 4.58$ & $\bigl[\begin{smallmatrix}9&2\\4&1\end{smallmatrix}\bigr]$ \\
$8$  & $16$ & $252$ & $6$ & $7$ &
  $\ell = 5.54$ & $\bigl[\begin{smallmatrix}15&7\\2&1\end{smallmatrix}\bigr]$ \\
\bottomrule
\end{tabular}
\end{center}
The terminal case $N = 23$ ($c = 3$) corresponds to
$Q^2 = \bigl[\begin{smallmatrix}2&1\\1&1\end{smallmatrix}\bigr]$,
the square of the standard Fibonacci matrix
$Q = \bigl[\begin{smallmatrix}1&1\\1&0\end{smallmatrix}\bigr]$;
the eigenvalues of $Q^2$ are $\varphi^2$ and $\varphi^{-2} = \xi^*(23)$.
$N = 23$ is the unique case with conductor $f = 1$:
the threshold $\xi^*$ is a unit in the maximal order $O_K$, not
merely in a suborder.
\end{proposition}

\begin{proof}
A $2 \times 2$ integer matrix with trace $c$ and determinant~$1$
lies in $\mathrm{SL}(2,\mathbb{Z})$; its characteristic polynomial
is $\lambda^2 - c\,\lambda + 1 = 0$, the same as the palindromic
threshold equation.  The existence of such a matrix for each
$c \in \mathbb{Z}_{\ge 3}$ is standard (every hyperbolic conjugacy
class contains a representative with the given trace).
The trace formula $c = 2 + 16/(N{-}7)$ follows from
$c = -B/A = 2(N{+}1)/(N{-}7)$ (Proposition~\ref{prop:alg-int})
by writing $2(N{+}1)/(N{-}7) = 2 + 16/(N{-}7)$.
The conductor $f$ satisfies $c^2 - 4 = f^2 D$ where
$D = \mathrm{sq\_free}(c^2 - 4)$.
Among the five traces $c \in \{16, 10, 6, 4, 3\}$,
the monic discriminants $c^2 - 4 = \{252, 96, 32, 12, 5\}$
have squarefree parts $\{7, 6, 2, 3, 5\}$ and
conductors $f = \{6, 4, 4, 2, 1\}$.
Only $c = 3$ ($N = 23$) gives $f = 1$,
since $5$ is the only squarefree value in the list.
\end{proof}

\begin{remark}[Automorphic correction on quotient surfaces]
\label{rmk:automorphic-correction}
On a quotient surface $\Gamma\backslash\mathbf{H}^2$, the
point vortex interaction involves the automorphic Green's function
$G_\Gamma(z,w) = \sum_{\gamma \in \Gamma} G(z, \gamma w)$.
Each non-trivial $\gamma$ contributes a correction to $C_1$
whose leading term is proportional to
$\mathrm{csch}^2(\ell_\gamma/2) = 4/\Delta_\gamma$,
where $\Delta_\gamma = \mathrm{tr}(\gamma)^2 - 4$ is the
discriminant of the geodesic
(since $\sinh(\ell/2) = \sqrt{\Delta}/2$).
The conductor decomposition
$\Delta = f^2 D$ gives
$\mathrm{csch}^2(\ell/2) = 4/(f^2 D)$:
the correction strength is the inverse discriminant,
factored through the conductor.
The palindromic structure of the threshold equation on
$\mathbf{H}^2$ reflects the conformal inversion
$\xi \leftrightarrow 1/\xi$; this symmetry is generically
\emph{broken} on quotient surfaces (it is preserved only
when $\Gamma$ is invariant under the inversion).
The stability thresholds on general quotients need not be
palindromic units.
\end{remark}

\begin{remark}[Topological sensitivity and the conductor]
\label{rmk:conductor-sensitivity}
The automorphic correction strength $4/\Delta = 4/(f^2 D)$
decreases rapidly with the conductor~$f$:
for $N = 23$ ($f = 1$, $\Delta = 5$) it is $0.80$, while
for $N = 8$ ($f = 6$, $\Delta = 252$) it is $0.016$.
The conductor thus measures the
\emph{topological sensitivity} of the stability threshold:
small $f$ (the golden ratio case) couples strongly to the
global topology of the surface, while large $f$ (the $N = 8$
case) is effectively insensitive.
The conductor is an intrinsic invariant of the $N \to (N{+}1)$
transition---it depends on the $\mathbb{Z}_N$ Casimir but not
on the surface.
\end{remark}

\paragraph{Higher-genus generalisation.}
The palindromic structure extends to vortex stability on
\emph{any} arithmetic hyperbolic surface.
If the Fuchsian group $\Gamma$ has trace field $K$ of degree $d$
over~$\mathbb{Q}$, the geodesic trace $B$ lies in~$K$, and
eliminating $B$ from $\xi^2 - B\xi + 1 = 0$ and the degree-$d$
minimal polynomial of $B$ yields a \emph{palindromic} polynomial
of degree $2d$ in $\xi$ over~$\mathbb{Q}$.
(Palindromic because $\xi \to 1/\xi$ fixes $B = \xi + 1/\xi$.)
Its roots are algebraic units of degree~$2d$.

\begin{proposition}[Higher-genus palindromic hierarchy]
\label{prop:higher-genus}
On an arithmetic surface $\Gamma\backslash\mathbf{H}^2$ with
trace field of degree $d$, the vortex stability threshold
satisfies a palindromic polynomial of degree $2d$
over~$\mathbb{Q}$.
For $d = 1$ (the modular surface): palindromic quadratics
and Pell units (the single-ring results of this paper).
For $d = 2$ (the Bolza surface, genus~$2$, trace field
$\mathbb{Q}(\sqrt{2})$): the systole has trace
$B = 2 + 2\sqrt{2}$
(\citealt{Buser1992}, \S13.3; the Bolza surface is the
genus-$2$ surface of maximal automorphism order~$48$, with
systole length $\ell = 2\,\mathrm{arccosh}(1{+}\sqrt{2})$),
giving the palindromic quartic
$\xi^4 - 4\xi^3 - 2\xi^2 - 4\xi + 1 = 0$
with threshold $\xi \approx 0.217$ in the quartic field
$\mathbb{Q}(\sqrt{2},\,\sqrt{2{+}2\sqrt{2}})$.
\end{proposition}

\begin{proof}
The substitution $u = \xi + 1/\xi$ maps $\xi^2 - u\xi + 1 = 0$
(palindromic quadratic) to a scalar equation $u = B$.
If $B$ has minimal polynomial $P(u) = 0$ of degree $d$
over~$\mathbb{Q}$, substituting $u = \xi + 1/\xi$ and
multiplying by $\xi^d$ gives $Q(\xi) = \xi^d\,P(\xi + 1/\xi)$,
a degree-$2d$ polynomial.
It is palindromic: $\xi^{2d}\,Q(1/\xi) = Q(\xi)$, because
$P(\xi + 1/\xi) = P(1/\xi + \xi)$.
The constant term of $Q$ equals $P(2)$ up to sign (from the
$\xi = 0$ limit), and the leading coefficient equals $1$
(from the $\xi \to \infty$ limit), so the roots are algebraic
units (product of all roots $= \pm 1$).
For the Bolza surface: $P(u) = u^2 - 4u - 4$
(minimal polynomial of $2 + 2\sqrt{2}$), giving
$Q(\xi) = \xi^4 - 4\xi^3 - 2\xi^2 - 4\xi + 1$, verified
palindromic with coefficients $[1, -4, -2, -4, 1]$.
\end{proof}

\paragraph{Arithmetic surfaces and quadratic reciprocity.}
On the congruence surface $\Gamma_0(p)\backslash\mathbf{H}^2$,
the geodesic of trace $B$ (and discriminant
$\Delta = B^2 - 4$) has a $\Gamma_0(p)$-representative if and
only if the Legendre symbol $(\Delta/p) \ne -1$.

\begin{proposition}[Quadratic reciprocity as a selection rule]
\label{prop:legendre-selection}
Of the five algebraic-integer stability thresholds, the number
visible on $\Gamma_0(p)\backslash\mathbf{H}^2$ depends on $p$:
\begin{center}
\begin{tabular}{ccccccc}
\toprule
$p$ & $N\!=\!23$ & $N\!=\!15$ & $N\!=\!11$ & $N\!=\!9$ & $N\!=\!8$
  & Count \\
& $\Delta\!=\!5$ & $\Delta\!=\!12$ & $\Delta\!=\!32$
  & $\Delta\!=\!96$ & $\Delta\!=\!252$ & \\
\midrule
$2$  & $\times$ & $\circ$ & $\circ$ & $\circ$ & $\circ$ & 4 \\
$3$  & $\times$ & $\circ$ & $\times$ & $\circ$ & $\circ$ & 3 \\
$5$  & $\circ$  & $\times$ & $\times$ & $\checkmark$ & $\times$ & 2 \\
$7$  & $\times$ & $\times$ & $\checkmark$ & $\times$ & $\circ$ & 2 \\
$11$ & $\checkmark$ & $\checkmark$ & $\times$ & $\times$ & $\times$ & 2 \\
$23$ & $\times$ & $\checkmark$ & $\checkmark$ & $\checkmark$ & $\times$ & 3 \\
$47$ & $\times$ & $\checkmark$ & $\checkmark$ & $\checkmark$ & $\checkmark$ & 4 \\
\bottomrule
\end{tabular}
\end{center}
($\checkmark$ = split, $\circ$ = ramified, $\times$ = inert.)
When $(\Delta/p) = +1$ (split), the geodesic lifts to
\emph{two} geodesics on $\Gamma_0(p)\backslash\mathbf{H}^2$,
doubling the automorphic correction.
When $(\Delta/p) = -1$ (inert), the geodesic is absent and
contributes \emph{zero} correction---the threshold on the quotient
matches the universal-cover value.
No prime $p < 200$ has all five thresholds simultaneously visible.
\end{proposition}

\begin{proof}
\emph{Step 1: conjugacy and quadratic forms.}
A hyperbolic $\gamma = \bigl[\begin{smallmatrix}a&b\\c&d
\end{smallmatrix}\bigr] \in \mathrm{SL}(2,\mathbb{Z})$
with trace $B = a{+}d$ has a conjugate in $\Gamma_0(p)$
iff some $\sigma\gamma\sigma^{-1}$ has lower-left entry
divisible by~$p$.
Conjugation by $\sigma = \bigl[\begin{smallmatrix}
\alpha&\beta\\\gamma'&\delta'\end{smallmatrix}\bigr]
\in \mathrm{SL}(2,\mathbb{Z})$ changes the lower-left
entry to $c' = c\alpha^2 + (d{-}a)\alpha\gamma' - b\gamma'^2$
(from the $(2,1)$-entry of $\sigma\gamma\sigma^{-1}$).
Hence $\gamma$ has a $\Gamma_0(p)$-conjugate iff the
binary quadratic form
$Q(x,y) = cx^2 + (d{-}a)xy - by^2$
represents a multiple of~$p$
(i.e., $Q(x,y) \equiv 0 \pmod{p}$ for some
$(x,y) \not\equiv (0,0) \pmod{p}$).

\emph{Step 2: discriminant and the Legendre symbol.}
The form $Q$ has discriminant
$(d{-}a)^2 + 4bc = (d{-}a)^2 + 4(ad{-}1) = (a{+}d)^2 - 4
= B^2 - 4 = \Delta$.
The equation $Q(x,y) \equiv 0 \pmod{p}$ with $y \not\equiv 0$
reduces (dividing by~$y^2$) to $cz^2 + (d{-}a)z - b \equiv 0$
where $z = x/y$, a quadratic in~$z$ with discriminant $\Delta$.
This has a solution mod~$p$ iff $\Delta$ is a quadratic residue
mod~$p$ (or $p \mid \Delta$):
precisely the condition $(\Delta/p) \ne -1$
(\citealt{Buser1992}, \S6; see also
\citealt{Katok1992}, Theorem~3.5.5 for the correspondence
between conjugacy classes in $\Gamma_0(N)$ and equivalence
classes of binary quadratic forms).

\emph{Step 3: trichotomy.}
When $(\Delta/p) = +1$ (split): $Q \equiv 0$ has two distinct roots
mod~$p$, giving \emph{two} $\Gamma_0(p)$-conjugacy classes in the
$\mathrm{SL}(2,\mathbb{Z})$ class of~$\gamma$---hence two geodesics
on $\Gamma_0(p)\backslash\mathbf{H}^2$.
When $p \mid \Delta$ (ramified): one double root, one geodesic.
When $(\Delta/p) = -1$ (inert): no root, no
$\Gamma_0(p)$-representative, no geodesic.

The table entries follow from evaluating $(\Delta/p)$ for
$\Delta \in \{5, 12, 32, 96, 252\}$ at each~$p$.
By the Chebotarev density theorem, each $\Delta$ is a QR for
asymptotically half of all primes; the five symbols are controlled
by four independent Legendre symbols $(2/p)$, $(3/p)$, $(5/p)$,
$(7/p)$ (since $96 = 32 \cdot 3$ gives $(\Delta_9/p) =
(\Delta_{11}/p)(\Delta_{15}/p)$), so the probability that all
five are simultaneously visible is $\sim (1/2)^4 = 1/16$,
predicting the first fully visible prime near
$p \approx 300$ (computed: $p = 311$).
\end{proof}

\noindent
The physical interpretation: on a hyperbolic lattice whose
global topology is that of $\Gamma_0(p)\backslash\mathbf{H}^2$,
only the \emph{visible} thresholds receive corrections from
the surface's geodesic spectrum.  The \emph{inert} thresholds
are identical to their universal-cover values---they are
\emph{topologically protected} by the Legendre symbol.
Different lattice topologies (different $p$) yield different
patterns of protected and unprotected thresholds.

\begin{corollary}[Characterization of fully visible primes]
\label{cor:fully-visible}
A prime $p > 7$ has all five thresholds simultaneously visible
on $\Gamma_0(p)\backslash\mathbf{H}^2$ if and only if
\[
  p \equiv r \pmod{840}
\]
for one of exactly $12$ residue classes $r$ out of
$\varphi(840) = 192$ coprime residues.
These $12$ classes are determined by the simultaneous conditions
\begin{alignat*}{2}
  (2/p) &\ne -1 &\quad&\Longleftrightarrow\quad p \equiv \pm 1 \pmod{8}, \\
  (3/p) &\ne -1 &&\Longleftrightarrow\quad p \equiv \pm 1 \pmod{12}, \\
  (5/p) &\ne -1 &&\Longleftrightarrow\quad p \equiv \pm 1 \pmod{5}, \\
  (7/p) &\ne -1 &&\Longleftrightarrow\quad
    \begin{cases} p \bmod 7 \in \{1,2,4\} & \text{if } p\equiv 1\!\!\pmod{4},\\
                  p \bmod 7 \in \{3,5,6\} & \text{if } p\equiv 3\!\!\pmod{4}.
    \end{cases}
\end{alignat*}
The density of fully visible primes is exactly $12/192 = 1/16$,
matching the independence prediction $(1/2)^4$ from the four
Legendre symbols.
The first fully visible prime is $p = 311$ $(\equiv 311 \bmod 840)$;
the first twenty are $311$, $479$, $719$, $839$, $1009$, $1129$,
$1151$, $1201$, $1319$, $1511$, $1559$, $1801$, $2351$, $2399$,
$2521$, $2689$, $2999$, $3049$, $3191$, $3359$.
\end{corollary}

\begin{proof}
The four conditions follow from the quadratic reciprocity law
applied to each prime factor of the discriminant set
$\{5, 12, 32, 96, 252\}$, as in the proof of
Proposition~\ref{prop:legendre-selection}.
The joint period is $\mathrm{lcm}(8, 12, 5, 28) = 840$
(the factor $28$ arises from the $(7/p)$ condition, which depends
on both $p \bmod 7$ and $p \bmod 4$, hence on $p \bmod 28$).
Each condition independently excludes half the coprime residues,
and the four conditions are algebraically independent
(the primes $\{2, 3, 5, 7\}$ generate independent
$\mathbb{Z}/2$-factors in
$(\mathbb{Z}/840\mathbb{Z})^\times$),
giving $192/2^4 = 12$ fully visible classes.
The listed primes are verified by direct Legendre symbol computation.
\end{proof}

\begin{theorem}[Universality of $N_{\mathrm{crit}} = 7$]
\label{thm:universality}
For any discrete subgroup $\Gamma \subset \mathrm{PSL}(2,\mathbb{R})$,
any $N \ge 3$, and any Fourier mode $m = 1, \ldots, N{-}1$:
the automorphic correction to the Havelock eigenvalue vanishes
at $\xi = 0$:
\[
  \delta C_1(\Gamma, N, m, \xi{=}0) = 0.
\]
In particular, $N_{\mathrm{crit}} = 7$ at $\xi = 0$ on every
quotient surface $\Gamma\backslash\mathbf{H}^2$, regardless of
the topology of the surface or the arithmeticity of~$\Gamma$.
\end{theorem}

\begin{proof}
At $\xi = 0$ ($r = 0$), all $N$ vortices sit at the origin
$z = 0$ of the Poincar\'e disk.  For any
$\gamma = \bigl[\begin{smallmatrix}a&b\\c&d\end{smallmatrix}\bigr]
\in \Gamma$:
$\gamma(0) = b/d$, a single point independent of the vortex
index~$k$.  The Fourier-weighted automorphic correction for
mode $m \ge 1$ is
\[
  \delta C_1(m)
  = \sum_{k=0}^{N-1} H_{rr}(0, b/d) \cdot \cos(2\pi mk/N)
  = H_{rr}(0, b/d) \cdot
  \underbrace{\sum_{k=0}^{N-1} \cos(2\pi mk/N)}_{= 0}
  = 0.
\]
The vanishing follows from $\mathbb{Z}_N$ character orthogonality
alone; no property of $\Gamma$ beyond discreteness is used.
Since $\delta C_1 = 0$ for every $\gamma \in \Gamma$ and every
mode, the Havelock eigenvalue at $\xi = 0$ on
$\Gamma\backslash\mathbf{H}^2$ equals its universal-cover value
$(N{-}1) - m(N{-}m)/2$ exactly.
\end{proof}

\begin{remark}[Topology as a curvature correction]
\label{rmk:topology-correction}
The theorem implies a hierarchy of corrections:
at $O(1)$ (the $\xi = 0$ limit), the Havelock eigenvalue is
universal---independent of the surface topology.
The global topology enters at $O(\xi)$: as the $N$-gon moves
away from the origin, the image vortices under $\Gamma$ spread
out (their positions $\gamma(z_k)$ become $k$-dependent), and
the Fourier projection becomes nonzero.
The $O(\xi)$ coefficient depends on $\Gamma$ and is controlled
by the Legendre symbol when $\Gamma$ is arithmetic
(Proposition~\ref{prop:legendre-selection}).
For non-arithmetic $\Gamma$, the coefficient exists but has no
known closed form.
\end{remark}

The coefficient $\gamma$ gives the curvature scale at which the polygon threshold is reached:
the $7\to 8$ jump occurs at $|K|\,r_E^2 = \xi^* = \gamma^{-1/2}$.
There is no perturbative formula $N_\mathrm{crit} \approx 7 + c\,|K|\,r_E^2$;
the threshold~\eqref{eq:xi_star} is a \emph{finite}, non-perturbative
condition, and $N_\mathrm{crit}$ is an integer that jumps discontinuously.

For any fixed ratio $R/a$ of ring radius to curvature radius, the
Gaussian curvature $K$ controls the threshold: $K > 0$ reduces
$N_{\mathrm{crit}}$ below~$7$, $K < 0$ raises it above~$7$, and
$K = 0$ gives exactly~$7$.  The number~$7$ is the stability threshold
at zero Gaussian curvature \emph{at the scale of the polygon ring}.
$K=0$ is the \emph{least stable point}: the $N=7$ ring is marginal
at $K=0$ and strictly stable for any $K<0$.

\begin{proposition}[Curvature-stability dichotomy]
\label{prop:curvature-dichotomy}
On a constant-curvature surface with Gaussian curvature $K$,
for a ring of $N$ vortices at any nonzero geodesic radius:
\[
  \mathrm{sgn}(N_{\mathrm{crit}} - 7) = -\mathrm{sgn}(K).
\]
That is: positive curvature destabilises
($N_{\mathrm{crit}} \le 7$), zero curvature is critical
($N_{\mathrm{crit}} = 7$), and negative curvature stabilises
($N_{\mathrm{crit}} \ge 7$), with strict inequality at any
nonzero ring radius.
\end{proposition}

\begin{proof}
On $\mathbf{S}^2$: $C_1 = (N{-}1)\cos\varphi < N{-}1$ for
$\varphi > 0$, so $C_1 < f_{\max}(N)$ is violated at smaller~$N$.
On $\mathbf{H}^2$: $C_1 = (N{-}1)(1{+}\xi^2)/(1{-}\xi)^2 > N{-}1$
for $\xi > 0$, so $C_1 \ge f_{\max}(N)$ holds at larger~$N$.
Both are monotone in their curvature parameter.
\end{proof}

The \emph{qualitative} classification
$\{N_{\mathrm{crit}} < 7,\; = 7,\; > 7\}$ depends only on
$\mathrm{sgn}(K)$, which for constant-curvature surfaces is a
topological invariant (it determines whether the universal cover
is $\mathbf{S}^2$, $\mathbb{R}^2$, or~$\mathbf{H}^2$).
In this sense the stability boundary $N = 7$ separates two
\emph{topological phases} of vortex polygon stability.
The \emph{quantitative} value $N_{\mathrm{crit}}(K, R)$ depends
on the dimensionless product $K R^2$ and is purely geometric.
\begin{proposition}[Morse index and spectral flow]
\label{prop:morse-index}
The Morse index of the $N$-gon critical point of the Thomson energy
on the constrained configuration space is
\[
  \mu(N) = \begin{cases} 0 & N \le 7, \\ N - 5 & N \ge 8. \end{cases}
\]
The unstable modes on the flat plane are exactly $m = 3, 4, \ldots, N{-}3$.
\end{proposition}

\begin{proof}
$\lambda_m < 0$ iff $m(N{-}m) > 2(N{-}1)$ iff
$(m - N/2)^2 < (N^2 - 8N + 8)/4$.  The left boundary of the
unstable band is $m_{\min} = \lceil N/2 - \sqrt{(N{-}4)^2 - 8}/2\,\rceil$.
For $N \ge 8$: $m_{\min} = 3$, because
$(N{-}6)^2 \le (N{-}4)^2 - 8$ iff $N \ge 7$
(giving $m_{\min} \le 3$) and
$(N{-}4)^2 > (N{-}4)^2 - 8$
(giving $m_{\min} > 2$).
By the palindromic symmetry $\lambda_m = \lambda_{N-m}$, the
upper boundary is $m_{\max} = N - 3$.
Hence $\mu(N) = (N{-}3) - 3 + 1 = N - 5$.
\end{proof}

\begin{corollary}[Why $N_{\mathrm{crit}} = 7$]
\label{cor:why-7}
$N_{\mathrm{crit}} = 7$ is the unique solution of $f(3, N) = C_1$
on the flat plane.
\end{corollary}
\begin{proof}
By Proposition~\ref{prop:morse-index}, the first unstable mode is
always $m = 3$.  Setting $f(3, N) = 3(N{-}3)/2 = N{-}1 = C_1$ gives
$N = 7$.
\end{proof}

The family $H_N(\xi) = C_1(\xi)\,I - F_N$ parametrised by
$\xi \in [0, 1)$ has a monotone spectral flow: since
$C_1(\xi) = (N{-}1)(1{+}\xi^2)/(1{-}\xi)^2$ satisfies
$\dd C_1/\dd\xi = (N{-}1)\cdot 2(1{+}\xi)/(1{-}\xi)^3 > 0$
for all $\xi \in [0,1)$, each eigenvalue
$\lambda_m(\xi) = C_1(\xi) - m(N{-}m)/2$ is strictly increasing
in~$\xi$.  Every eigenvalue that is negative at $\xi = 0$
(flat plane) crosses zero exactly once as $\xi \to 1$ (infinite
curvature).  The total spectral flow therefore equals the
flat-plane Morse index:
$\mathrm{SF}(H_N, 0 \to 1) = \mu(N)$.  This is \emph{proven}
(it follows from Proposition~\ref{prop:morse-index} together with
the strict monotonicity of $C_1(\xi)$, which is an explicit
calculation---not the $N_{\mathrm{crit}}$ monotonicity conjecture
of Section~\ref{sec:curvature-thresholds}, which concerns the
integer-valued function $N_{\mathrm{crit}}(R/a)$, a different
quantity).

\begin{remark}[Robustness of $\mu(N)$ and the $N = 7$ exception]
\label{rmk:topological-protection}
For $N \ne 7$, all eigenvalues $\lambda_m(0)$ are nonzero, so the
spectral flow $\mathrm{SF} = \mu(N)$ is determined entirely by the
boundary data (the Havelock spectrum at $\xi = 0$ and the
divergence $C_1 \to \infty$ at $\xi \to 1^-$).
By standard Fredholm index stability
(\citealt{Kato1995}, Theorem~IV.5.17),
$\mu(N)$ is insensitive to any bounded perturbation $V(\xi)$
with $V(0) = 0$ that couples Fourier modes in the interior---even
perturbations that break the $\mathbb{Z}_N$ diagonal structure
and create avoided crossings cannot change the net spectral flow
(\texttt{proofs/topological\_protection.py} verifies this for
$N = 8$--$10$ with coupling strengths up to $\|V\| = 100$).
At $N = 7$, the kernel of $H_7(0)$ (modes $m = 3, 4$ with
$\lambda_m = 0$) breaks this robustness:
sufficiently strong mode coupling can push the zero eigenvalues
negative, changing the spectral flow.
The heptagonal ring is the unique case where the Morse index is
structurally fragile.
\end{remark}

\begin{theorem}[Algebraic origin of $N_{\mathrm{crit}} = 7$]
\label{thm:why-7-algebraic}
$N_{\mathrm{crit}} = 7$ is the unique positive integer $N \ge 3$ at
which the $\mathbb{Z}_N$ Casimir maximum equals the logarithmic $C_1$:
\begin{equation}\label{eq:n-crit-factorization}
  \max_m f(m,N) = C_1
  \;\Longleftrightarrow\;
  N^2 - 8N + 7 = (N-1)(N-7) = 0.
\end{equation}
\end{theorem}

\begin{proof}
For odd $N$ (the binding case): $\max_m f = (N^2{-}1)/8$
and $C_1 = N{-}1$.  Setting $(N^2{-}1)/8 = N{-}1$ and
multiplying by~$8$: $N^2 - 1 = 8(N{-}1)$, i.e.,
$N^2 - 8N + 7 = 0$.  This factors as $(N{-}1)(N{-}7) = 0$,
giving $N = 7$ (the root $N = 1$ is below the polygon range).
For even $N$: $\max_m f = N^2/8$; setting $N^2/8 = N{-}1$
gives $N^2 - 8N + 8 = 0$ with roots $4 \pm 2\sqrt{2}$
(irrational, so no even transition exists).
\end{proof}

\noindent
The coefficient $8 = 2 \cdot 4$ in~\eqref{eq:n-crit-factorization}
arises from two intrinsic features of $\mathbb{Z}_N$:
the Havelock normalization $f(m,N) = m(N{-}m)/\mathbf{2}$
and the critical-mode evaluation
$m = (N{-}1)/2$ giving $m(N{-}m) = (N^2{-}1)/\mathbf{4}$.
The factorization $(N{-}1)(N{-}7)$ is therefore purely
representation-theoretic: no differential geometry, no index
theorems, no PDEs---just the intersection of the quadratic
Casimir with the linear $C_1$ from dilation invariance.
The number~$7$ is an algebraic shadow of the $\mathbb{Z}_N$
representation theory combined with the uniqueness of the
logarithmic interaction.

\paragraph{Physical realization.}
The hyperbolic plane is no longer a purely mathematical object.
\citet{Kollar2019} demonstrated that microwave resonators arranged in
a tiling of the Poincar\'e disk implement the hyperbolic lattice
Hamiltonian in circuit quantum electrodynamics, with negative
effective curvature controlled by the lattice geometry.
\citet{Lenggenhager2022} subsequently implemented hyperbolic space
on a classical circuit board, showing that band structure and
topological properties native to $\mathbf{H}^2$ are experimentally
accessible.  These platforms realize, in a controlled laboratory
setting, the geometry that Section~\ref{sec:hyperbolic} analyzes.

The stability prediction translates directly into an experimental
observable: \emph{a ring of $N$ vortex-like excitations on a
hyperbolic lattice with ring-to-curvature ratio $R/a \approx 1$
should be linearly stable up to $N = 12$, without requiring any
central defect}.  In flat geometry, a 12-vortex ring is linearly
unstable at every curvature-free configuration and can only be
stabilized by a sufficiently strong central vortex
(the Thomson--Havelock criterion requires a central circulation
$\kappa_0 / \kappa \ge |\kappa_{\mathrm{crit}}(12)|$, which is large).
The hyperbolic curvature replaces the stabilizing role of the
central vortex: the geometry itself provides the restoring force.
A predicted $N = 12$ stable ring with no center is impossible in
flat geometry and would be a direct signature of hyperbolic curvature
acting on vortex dynamics.

\paragraph{Hyperbolic lattices as an isolated test of Thomson stability.}
In the planetary context, the three selection constraints---Rossby
stationarity, Thomson stability, and geometric packing---are all
active simultaneously and calibrated to the same observations;
disentangling their individual contributions requires model-dependent
assumptions.  A hyperbolic lattice experiment decouples them:
\begin{itemize}
  \item There is \emph{no} Rossby wave mechanism (no planetary
    rotation, no $\beta$-effect).
  \item Thomson stability is directly testable by counting stable
    vortex excitations as $R/a$ is varied.
  \item Geometric packing is controllable by engineering the
    impurity spacing independently of the curvature.
\end{itemize}
This makes hyperbolic lattices the cleanest possible test of the
Thomson constraint \emph{in isolation}.

Moreover, the algebraic-integer thresholds
(Proposition~\ref{prop:alg-int}) are \emph{exact, parameter-free}
predictions:
\begin{center}
\begin{tabular}{cccl}
\toprule
Transition & $\xi^*$ (exact) & $R/a$ & Prediction \\
\midrule
$7 \to 8$   & $8 - 3\sqrt{7} = 0.0627\ldots$ & $0.51$ &
  8-gon stabilises \\
$10 \to 11$ & $3 - 2\sqrt{2} = 0.1716\ldots$ & $0.88$ &
  11-gon stabilises \\
$14 \to 15$ & $2 - \sqrt{3}\phantom{0} = 0.2679\ldots$ & $1.15$ &
  15-gon stabilises \\
$22 \to 23$ & $(3{-}\sqrt{5})/2 = 0.3820\ldots$ & $1.44$ &
  23-gon stabilises \\
\bottomrule
\end{tabular}
\end{center}
The threshold values are algebraic numbers (Pell units and golden
ratio powers), not fitted parameters.  Observing any of these
transitions would constitute a direct experimental verification
of the palindromic--Pell--conformal-inversion structure identified
in Section~\ref{sec:algebraic-structure}.
The physical conditions required to realize the hyperbolic vortex
interaction in a circuit-QED or BEC platform are discussed in
companion Paper~II alongside the BEC predictions.

%% ====================================================================

\paragraph{Open problems.}
\begin{enumerate}
  \item \emph{Arithmetic surfaces and quadratic reciprocity.}
    On the quotient $\Gamma_0(p)\backslash\mathbf{H}^2$, the
    geodesic of trace $B$ has a $\Gamma_0(p)$-representative iff
    the Legendre symbol $(\Delta/p) \ne -1$ where $\Delta = B^2 - 4$.
    Which algebraic-integer thresholds are \emph{visible} on which
    arithmetic surface is therefore controlled by quadratic
    reciprocity---a selection rule from algebraic number theory
    acting on vortex stability.  Making this quantitative
    (computing the automorphic correction to $C_1$ on specific
    $\Gamma_0(p)$ quotients) is an open problem.
  \item \emph{Non-constant curvature.}
    On an oblate spheroid, the Gaussian curvature varies with
    latitude, breaking the $\mathbb{Z}_N$ decomposition and
    coupling Fourier modes.
    Paper~III shows this effect is critical for Saturn's hexagon
    ($10\%$ oblateness flips the sign of $\lambda_3$)
    but negligible for Neptune/Uranus.
    A full stability analysis using the spheroidal Green's function
    (oblate spheroidal harmonics) would extend the framework
    beyond constant curvature.
  \item \emph{Experimental verification on hyperbolic lattices.}
    The five algebraic-integer thresholds are parameter-free
    predictions for circuit-QED or topoelectric platforms
    implementing hyperbolic tilings.
    The most accessible test: a ring of $N = 12$ excitations
    at $R/a \approx 1$ should be stable---impossible on any
    flat-geometry platform.
\end{enumerate}

%% ====================================================================

\bibliographystyle{plainnat}
\bibliography{../shared/refs}

%% ====================================================================
%% ADDENDUM: The Bolza modular form
%% ====================================================================

\newpage
\appendix
\section{The Bolza modular form}
\label{sec:bolza-modular}

This addendum identifies the weight-$1$ modular form whose
Hecke eigenvalues encode the splitting pattern of the Bolza
palindromic quartic from Proposition~\ref{prop:higher-genus}.
It is the first non-abelian case in the palindromic hierarchy
and provides a concrete connection between vortex stability and
automorphic forms via the Langlands--Tunnell correspondence.

\subsection{The palindromic quartic and its Galois group}

The Bolza surface (the genus-$2$ surface of maximal automorphism
order~$48$; \citealt{Buser1992}, \S13.3) has systole trace
$B = 2 + 2\sqrt{2}$, which satisfies $u^2 - 4u - 4 = 0$.
By the palindromic construction of
Proposition~\ref{prop:higher-genus}, the stability threshold
satisfies the palindromic quartic
\begin{equation}\label{eq:bolza-quartic}
  \xi^4 - 4\xi^3 - 2\xi^2 - 4\xi + 1 = 0,
\end{equation}
with coefficients $[1, -4, -2, -4, 1]$ (palindromic).
The splitting field is $K = \mathbb{Q}(\sqrt{2},\,
\sqrt{2{+}2\sqrt{2}})$, generated by a root $\alpha$
of $\alpha^4 - 4\alpha^2 - 4 = 0$.

\begin{proposition}[Galois group of the Bolza quartic]
\label{prop:bolza-galois}
$\mathrm{Gal}(K/\mathbb{Q}) \cong D_4$ (dihedral of order~$8$).
The field discriminant is $\mathrm{disc}(K/\mathbb{Q}) = -2^{10}$.
Only the prime $p = 2$ ramifies.
\end{proposition}

\begin{proof}
The resolvent cubic of the depressed quartic $y^4 + by^2 + d$
(obtained from $\alpha^4 - 4\alpha^2 - 4$) is
$t^3 + 4t^2 + 16t = t(t^2 + 4t + 16)$.
The quadratic factor has discriminant $16 - 64 = -48 < 0$
(no real roots beyond $t = 0$), so the resolvent has exactly
one rational root.
The polynomial discriminant
$\mathrm{disc}(\alpha^4 - 4\alpha^2 - 4) = -2^{16}$ is
negative, so $\mathrm{Gal} \not\subset A_4$.
A quartic with one rational resolvent root and Galois group
$\not\subset A_4$ has $\mathrm{Gal} \cong D_4$
(\citealt{Katok1992}).
The field discriminant $\mathrm{disc}(K/\mathbb{Q}) = -2^{10}$
is computed from the integral basis
$\{1,\, \alpha,\, \alpha^2/2,\, (\alpha + \alpha^3)/4\}$
(verified in Mathematica;
$[\mathcal{O}_K : \mathbb{Z}[\alpha]]^2
= 2^{16}/2^{10} = 2^6$, so $[\mathcal{O}_K : \mathbb{Z}[\alpha]] = 8$).
\end{proof}

\subsection{The Artin representation and Langlands--Tunnell}

The group $D_4$ has five irreducible representations:
four one-dimensional ($\mathrm{triv}$, $\mathrm{sign}$,
$\chi_1$, $\chi_2$) and one two-dimensional ($\rho$).
The character table is:
\begin{center}
\begin{tabular}{lccccc}
\toprule
& $1$ & $r^2$ & $r$ & $s$ & $rs$ \\
\midrule
$\mathrm{triv}$ & $1$ & $1$ & $1$ & $1$ & $1$ \\
$\mathrm{sign}$ & $1$ & $1$ & $-1$ & $1$ & $-1$ \\
$\chi_1$ & $1$ & $1$ & $1$ & $-1$ & $-1$ \\
$\chi_2$ & $1$ & $1$ & $-1$ & $-1$ & $1$ \\
$\rho$   & $2$ & $-2$ & $0$ & $0$ & $0$ \\
\bottomrule
\end{tabular}
\end{center}
The Artin $L$-function $L(s, \rho)$ has conductor
$N = 2^7 = 128$ (the Artin conductor divides
$|\mathrm{disc}(K/\mathbb{Q})| = 2^{10}$ and is determined
by the higher ramification filtration at $p = 2$;
the LMFDB lists the conductor as~$128$ for field
\texttt{4.2.1024.1}).

Since $D_4$ is solvable, the Langlands--Tunnell theorem
\citep{Tunnell1981} guarantees that $L(s,\rho) = L(s,f)$
for a weight-$1$ modular form $f$ on $\Gamma_1(256)$.

\subsection{Identification: the form 256.1.c.a}

\begin{theorem}[Bolza modular form]
\label{thm:bolza-modular}
The Artin $L$-function of the two-dimensional representation
$\rho$ of $D_4 = \mathrm{Gal}(K/\mathbb{Q})$ corresponds to
the weight-$1$ newform
\[
  f = \eta(8z)\,\eta(16z)
  = q - q^9 - 2q^{17} + q^{25} + 2q^{41} + q^{49}
    - 2q^{73} + q^{81} - 2q^{89} - 2q^{97} + O(q^{100})
\]
of level $128$, character $\chi_{128.d}$ (order~$2$),
with rational Fourier coefficients.
This form is listed as \texttt{128.1.d.a} in the LMFDB,
and its Stark unit satisfies the palindromic quartic
$\xi^4 - 4\xi^3 - 2\xi^2 - 4\xi + 1 = 0$---exactly
equation~\eqref{eq:bolza-quartic}.
\end{theorem}

\begin{proof}
The Hecke eigenvalue $a_p = \mathrm{tr}(\rho(\mathrm{Frob}_p))$
is determined by the splitting of $p$ in $K$:
\begin{itemize}
  \item $p \equiv 3, 5 \pmod{8}$ ($p$ inert in
    $\mathbb{Q}(\sqrt{2})$): $\mathrm{Frob}_p$ has order~$4$
    in $D_4$, so $a_p = 0$.
  \item $p \equiv \pm 1 \pmod{8}$ ($p$ splits in
    $\mathbb{Q}(\sqrt{2})$):
    \begin{itemize}
      \item $\alpha^4 - 4\alpha^2 - 4$ has $4$ roots mod $p$
        $\Rightarrow$ $p$ splits completely in $K$,
        $\mathrm{Frob}_p = 1$, $a_p = 2$;
      \item $2$ roots mod $p$ $\Rightarrow$ partial splitting,
        $\mathrm{Frob}_p$ is a reflection, $a_p = 0$;
      \item $0$ roots mod $p$ $\Rightarrow$ $\mathrm{Frob}_p = r^2$
        (central element), $a_p = -2$.
    \end{itemize}
\end{itemize}
\begin{center}
\begin{tabular}{ccccl}
\toprule
$p$ & $p \bmod 8$ & Roots mod $p$ & $a_p$ & $D_4$ class \\
\midrule
$7$  & $7$ & $2$ & $0$ & reflection \\
$17$ & $1$ & $0$ & $-2$ & $r^2$ \\
$23$ & $7$ & $2$ & $0$ & reflection \\
$41$ & $1$ & $4$ & $\mathbf{+2}$ & identity \\
$73$ & $1$ & $0$ & $-2$ & $r^2$ \\
$89$ & $1$ & $0$ & $-2$ & $r^2$ \\
$97$ & $1$ & $0$ & $-2$ & $r^2$ \\
\bottomrule
\end{tabular}
\end{center}
The $q$-expansion of the eta product
$\eta(8z)\eta(16z)$, computed independently
(LMFDB, \texttt{128.1.d.a}), matches these values at every
prime $p < 100$.
The Chebotarev density theorem predicts the asymptotic
distribution $a_p = 2$ for $1/8$ of primes, $a_p = -2$ for
$1/8$, and $a_p = 0$ for $3/4$---verified numerically
for $p < 1000$.
\end{proof}

\begin{remark}[Physical content and spectral bound]
\label{rmk:bolza-physics}
The Hecke eigenvalue $a_p$ of the form
$f = \texttt{128.1.d.a}$ determines the automorphic correction
to vortex stability on the Bolza surface at each prime~$p$:
$a_p = 2$ (complete splitting, doubled correction),
$a_p = 0$ (partial or inert, no net correction),
$a_p = -2$ (central inertia, reversed correction).
This is the non-abelian analogue of the Legendre symbol
selection rule in Proposition~\ref{prop:legendre-selection}:
the genus-$0$ case uses Dirichlet characters (abelian,
class field theory), while the genus-$2$ Bolza case requires
a genuine weight-$1$ modular form (non-abelian, Langlands--Tunnell).

The palindromic quartic~\eqref{eq:bolza-quartic} is derived on
$\mathbf{H}^2$ (the universal cover), not on the compact Bolza
surface.  The spectral expansion of the automorphic correction
$\delta C_1 = \sum_{n \ge 1} A_n/\lambda_n$ (summing over
Bolza Laplacian eigenfunctions) is bounded by the
\emph{trivial-representation} eigenvalues: only eigenfunctions
invariant under the full automorphism group (order~$48$) are
nonzero at the centre of the fundamental domain
(confirmed by numerical eigenfunction computation).
The first trivial-representation eigenvalue of the Bolza surface
is $\lambda_{\mathrm{triv}} \approx 23.08$
(\citealt{Strohmaier2013}, multiplicity-$1$ entry), giving
$|\delta C_1| \lesssim 2/(4\pi\,\lambda_{\mathrm{triv}})
\approx 0.16$.
Since the binding eigenvalue at $N_{\mathrm{crit}} = 12$
is $\lambda_{\mathrm{bind}} = 0.78$, the correction satisfies
$|\delta C_1| < \lambda_{\mathrm{bind}}$:
the compact-surface correction does not change $N_{\mathrm{crit}}$,
and the $\mathbf{H}^2$ palindromic quartic remains the correct
threshold polynomial on the Bolza surface.
\end{remark}

\subsection{The Fibonacci threshold tower}

The Bolza modular form corresponds to a non-congruence surface.
On \emph{congruence} surfaces $\Gamma(p)\backslash\mathbf{H}^2$,
a different structure emerges: the Fibonacci matrix controls the
threshold through the Pisano period.

\begin{theorem}[Fibonacci systole of $\Gamma(p)\backslash\mathbf{H}^2$]
\label{thm:fibonacci-systole}
Let $p$ be an odd prime with entry point
$\alpha(p) = \min\{n > 0 : p \mid F_n\}$.
Define $n = \alpha(p)$ if $\alpha(p)$ is even,
$n = 2\alpha(p)$ if $\alpha(p)$ is odd.
Then $\Gamma(p)$ contains an element of trace $\pm L_n$
(the $n$-th Lucas number) in the
$\mathrm{SL}(2,\mathbb{Z})$-conjugacy class of the
Fibonacci matrix, and the corresponding stability threshold
satisfies the palindromic quadratic
\[
  \xi^2 - L_n\,\xi + 1 = 0,
\]
with discriminant $L_n^2 - 4 = 5\,F_n^2$,
conductor $F_n$, and field $\mathbb{Q}(\sqrt{5})$.
\end{theorem}

\begin{proof}
The Fibonacci matrix $Q = \bigl[\begin{smallmatrix}1&1\\1&0
\end{smallmatrix}\bigr]$ has $Q^n =
\bigl[\begin{smallmatrix}F_{n+1}&F_n\\F_n&F_{n-1}
\end{smallmatrix}\bigr]$, $\mathrm{tr}(Q^n) = L_n$,
$\det(Q^n) = (-1)^n$.
Since $n$ is even, $\det = 1$.
By the definition of the entry point, $p \mid F_{\alpha(p)}$.
For even $\alpha$: $p \mid F_n$, and since
$F_{n \pm 1} \equiv \pm 1 \pmod{p}$ (from the Pisano periodicity),
either $Q^n \equiv I$ or $Q^n \equiv -I \pmod{p}$.
In the latter case, $(-I)Q^n \equiv I \pmod{p}$, giving
an element of $\Gamma(p)$ with trace $-L_n$.
For odd $\alpha$: $F_{2\alpha} = F_\alpha L_\alpha \equiv 0 \pmod{p}$,
and $Q^{2\alpha}$ is handled by the same argument.
The palindromic property, discriminant identity
$L_n^2 - 4 = 5F_n^2$ (a standard Fibonacci relation),
and field $\mathbb{Q}(\sqrt{5})$ follow as in
Proposition~\ref{prop:geodesic-correspondence}.
\end{proof}

\noindent
The theorem connects three classical sequences:
\begin{center}
\begin{tabular}{cccccl}
\toprule
$p$ & $\alpha(p)$ & $n$ & $|L_n|$ & $F_n$ & Conductor \\
\midrule
$3$ & $4$ & $4$ & $7$ & $3$ & $3$ \\
$7$ & $8$ & $8$ & $47$ & $21$ & $21$ \\
$11$ & $10$ & $10$ & $123$ & $55$ & $55$ \\
$29$ & $14$ & $14$ & $843$ & $377$ & $377$ \\
$47$ & $16$ & $16$ & $2207$ & $987$ & $987$ \\
$19$ & $18$ & $18$ & $5778$ & $2584$ & $2584$ \\
\bottomrule
\end{tabular}
\end{center}
All thresholds lie in $\mathbb{Q}(\sqrt{5})$ because
$\sqrt{L_n^2 - 4} = F_n\sqrt{5}$, and the conductor
is the Fibonacci number $F_n$.
The Pisano period $\pi(p)$, previously a purely number-theoretic
quantity (the period of $F_n \bmod p$), acquires geometric meaning:
it is the geodesic winding number of the Fibonacci geodesic
on $\Gamma(p)\backslash\mathbf{H}^2$.

\begin{remark}[The Takeuchi obstruction]
\label{rmk:takeuchi}
All $85$ arithmetic triangle groups in the Takeuchi
classification have trace fields contained in cyclotomic fields.
By the Kronecker--Weber theorem, their Galois groups are
abelian, so the palindromic polynomials from triangle groups
always have \emph{solvable} Galois groups.
In particular, the Langlands--Tunnell theorem applies to every
triangle-group surface.
The Bolza surface ($D_4$, non-abelian but solvable) is the
simplest case where a genuine automorphic form appears.
Non-solvable Galois groups (requiring the full Langlands
correspondence) would arise from quaternion algebras over
totally real number fields with non-abelian Galois group---a
direction requiring the arithmetic of non-congruence Fuchsian
groups.
\end{remark}

\begin{remark}[The stability fingerprint of a prime]
\label{rmk:fingerprint}
Each prime $p$ defines a \emph{stability fingerprint}: a function
from vortex polygon numbers $N$ to the set
$\{\text{stable, unstable, marginal}\}$ on the surface
$\Gamma(p)\backslash\mathbf{H}^2$.
For the five algebraic-integer thresholds of the modular surface
(Proposition~\ref{prop:geodesic-correspondence}), the fingerprint
is determined by four Legendre symbols $(2/p)$, $(3/p)$, $(5/p)$,
$(7/p)$ (Corollary~\ref{cor:fully-visible}), giving $2^4 = 16$
fingerprint classes with period $840$.
Adding the Fibonacci entry point $\alpha(p)$
(Theorem~\ref{thm:fibonacci-systole}) refines the $(5/p) = +1$
class into a family parameterised by $\alpha(p)$, with
conductor $F_{\alpha(p)}$.
Adding the Bolza $D_4$ threshold
(Theorem~\ref{thm:bolza-modular}) contributes a Hecke eigenvalue
$a_p \in \{-2, 0, 2\}$ from the weight-$1$ form $256.1.c.a$.

The full stability fingerprint of a prime is therefore a
composite invariant:
\begin{center}
\begin{tabular}{lll}
\toprule
Component & Type & Source \\
\midrule
$(2/p),\,(3/p),\,(5/p),\,(7/p)$ &
  Legendre symbols (abelian) &
  Prop.~\ref{prop:legendre-selection} \\
$\alpha(p)$ &
  Fibonacci entry point (arithmetic) &
  Thm.~\ref{thm:fibonacci-systole} \\
$a_p$ &
  $D_4$ Hecke eigenvalue (automorphic) &
  Thm.~\ref{thm:bolza-modular} \\
\bottomrule
\end{tabular}
\end{center}
This composite is computable for any given prime but has
statistical properties depending on deep conjectures: the
generalised Riemann hypothesis for the entry point distribution,
the Sato--Tate conjecture (in its $D_4$ form) for the Hecke
eigenvalue distribution.
The vortex stability problem organises these number-theoretic
invariants into a single coherent structure---the stability table
of an $N$-gon on a specific arithmetic surface.

Such correlations exist.  Among primes with $(2/p) = +1$
(the only primes where $a_p$ can be nonzero):
\begin{center}
\begin{tabular}{lccc}
\toprule
$\alpha(p)$ parity & $a_p = -2$ & $a_p = 0$ & $a_p = 2$ \\
\midrule
Even ($n = 149$ primes $< 3000$) & $16\%$ & $73\%$ & $11\%$ \\
Odd ($n = 61$ primes $< 3000$) & $48\%$ & $\mathbf{0\%}$ & $52\%$ \\
\bottomrule
\end{tabular}
\end{center}
When $\alpha(p)$ is odd, $a_p = 0$ \emph{never occurs}---an
exact constraint, not a statistical trend.

\begin{proposition}[Fibonacci--Bolza complementarity]
\label{prop:complementarity}
Let $p > 5$ be prime with $(2/p) = +1$.
\begin{enumerate}
  \item If\, $p \equiv 3 \pmod{4}$, then $\alpha(p)$ is even.
  \item If\, $p \equiv 1 \pmod{4}$ and $(5/p) = -1$, then
    $\alpha(p)$ is odd.
  \item If\, $\alpha(p)$ is odd and $(2/p) = +1$, then
    $a_p \in \{+2, -2\}$ (the Bolza Hecke eigenvalue is nonzero).
\end{enumerate}
In particular, for $p \equiv 17$ or $33 \pmod{40}$:
the golden ratio threshold is invisible $((5/p) = -1)$ and
the Bolza threshold is maximally active $(a_p \ne 0)$.
\end{proposition}

\begin{proof}
(i)~When $(-1/p) = -1$, the order of $-\varphi^2$ in
$\mathbb{F}_{p^2}^*$ divides $p + 1$ (even), and
$\alpha(p)$ is the order of $-\varphi^2$ modulo~$p$
in the appropriate residue ring, hence even.
(This is the standard parity result for the Fibonacci
entry point; see~\citealt{Ribenboim1996}, Theorem~2.)

(ii)~When $(-1/p) = +1$ and $(5/p) = -1$:
$\varphi \notin \mathbb{F}_p$ but $\varphi \in \mathbb{F}_{p^2}$.
The entry point $\alpha(p)$ divides $p + 1$
(since $\varphi^{p+1} = \mathrm{Norm}(\varphi) = -1$
in $\mathbb{F}_{p^2}/\mathbb{F}_p$, giving
$(-\varphi^2)^{p+1} = \varphi^{2(p+1)} = (-1)^2 = 1$).
Write $p + 1 = 2^s \cdot t$ with $t$ odd.
Since $(-1/p) = +1$ and $(5/p) = -1$, the order of
$-\varphi^2$ modulo~$p$ has odd part equal to $t$
(the full odd part of $p + 1$), so $\alpha(p)$ is odd.

(iii)~Odd $\alpha$ means $Q^\alpha$ has
$\det = (-1)^\alpha = -1$ and trace $L_\alpha$.
The element $(-I)\cdot Q^\alpha \in \Gamma(p)$ has trace
$-L_\alpha$ and $\det = 1$.
The eigenvalues of $Q^\alpha$ are $\varphi^\alpha$ and
$(-1/\varphi)^\alpha = -\varphi^{-\alpha}$ (note the sign).
Reducing modulo a prime $\mathfrak{P}$ of $K$ above~$p$:
the eigenvalue ratio is $-\varphi^{2\alpha}$, which is
$\equiv -1 \pmod{\mathfrak{P}}$ (from part~(ii):
$\varphi^{2\alpha} \equiv 1$).
An element with eigenvalue ratio $-1$ modulo $\mathfrak{P}$
acts as $-I$ on the residue field, placing
$\mathrm{Frob}_p$ in the centre $\{1, r^2\}$ of $D_4$
(where $\mathrm{tr}(\rho) = \pm 2$), not in a reflection
class (where $\mathrm{tr}(\rho) = 0$).

Verified for all primes $p < 10{,}000$ with zero violations.
\end{proof}

\begin{proposition}[Twin prime shared divisibility]
\label{prop:twin-fibonacci}
Let $(p, p{+}2)$ be a twin prime pair with $p \equiv 2 \pmod{5}$
(equivalently, $(5/p) = -1$ and $(5/(p{+}2)) = +1$).
Then both Fibonacci entry points divide $p + 1$:
\[
  \alpha(p) \mid (p{+}1)
  \qquad\text{and}\qquad
  \alpha(p{+}2) \mid (p{+}1).
\]
This shared constraint is unique to twin primes: for any prime
gap $g > 2$, the condition
$g = (5/(p{+}g)) - (5/p)$ has $|g| \le 2$ (since both
Legendre symbols are $\pm 1$), so no larger gap can produce it.
\end{proposition}

\begin{proof}
Since $(5/p) = -1$: $\alpha(p) \mid p - (5/p) = p + 1$.
Since $(5/(p{+}2)) = +1$: $\alpha(p{+}2) \mid (p{+}2) - 1 = p + 1$.
For the uniqueness: the shared constraint requires
$p - (5/p) = (p{+}g) - (5/(p{+}g))$, i.e.,
$g = (5/(p{+}g)) - (5/p)$.  Since Legendre symbols take values
$\pm 1$, $|g| \le 2$.  For $g = 2$: achievable when $(5/p) = -1$
and $(5/(p{+}2)) = +1$, which is $p \equiv 2 \pmod{5}$.
For $g = 4, 6, \ldots$: impossible.
\end{proof}

\noindent
The stability fingerprints of twin primes with $p \equiv 2 \pmod{5}$
are \emph{correlated} through the factorisation of $p + 1$:
both Fibonacci thresholds are controlled by the divisors of a
single integer.
Combined with Proposition~\ref{prop:complementarity}(ii), if also
$p \equiv 1 \pmod{4}$ (i.e., $p \equiv 17$ or $33 \pmod{40}$):
$\alpha(p)$ is odd and $a_p \in \{\pm 2\}$ (Bolza maximally active),
while $\alpha(p{+}2)$ divides the same $p + 1$ but its parity is
unconstrained---the twin pair has symmetric Fibonacci fingerprints
but potentially asymmetric Bolza fingerprints.
\end{remark}

\end{document}
